\documentclass[11pt]{article}

    \usepackage[breakable]{tcolorbox}
    \usepackage{parskip} % Stop auto-indenting (to mimic markdown behaviour)
    

    % Basic figure setup, for now with no caption control since it's done
    % automatically by Pandoc (which extracts ![](path) syntax from Markdown).
    \usepackage{graphicx}
    % Keep aspect ratio if custom image width or height is specified
    \setkeys{Gin}{keepaspectratio}
    % Maintain compatibility with old templates. Remove in nbconvert 6.0
    \let\Oldincludegraphics\includegraphics
    % Ensure that by default, figures have no caption (until we provide a
    % proper Figure object with a Caption API and a way to capture that
    % in the conversion process - todo).
    \usepackage{caption}
    \DeclareCaptionFormat{nocaption}{}
    \captionsetup{format=nocaption,aboveskip=0pt,belowskip=0pt}

    \usepackage{float}
    \floatplacement{figure}{H} % forces figures to be placed at the correct location
    \usepackage{xcolor} % Allow colors to be defined
    \usepackage{enumerate} % Needed for markdown enumerations to work
    \usepackage{geometry} % Used to adjust the document margins
    \usepackage{amsmath} % Equations
    \usepackage{amssymb} % Equations
    \usepackage{textcomp} % defines textquotesingle
    % Hack from http://tex.stackexchange.com/a/47451/13684:
    \AtBeginDocument{%
        \def\PYZsq{\textquotesingle}% Upright quotes in Pygmentized code
    }
    \usepackage{upquote} % Upright quotes for verbatim code
    \usepackage{eurosym} % defines \euro

    \usepackage{iftex}
    \ifPDFTeX
        \usepackage[T1]{fontenc}
        \IfFileExists{alphabeta.sty}{
              \usepackage{alphabeta}
          }{
              \usepackage[mathletters]{ucs}
              \usepackage[utf8x]{inputenc}
          }
    \else
        \usepackage{fontspec}
        \usepackage{unicode-math}
    \fi

    \usepackage{fancyvrb} % verbatim replacement that allows latex
    \usepackage{grffile} % extends the file name processing of package graphics
                         % to support a larger range
    \makeatletter % fix for old versions of grffile with XeLaTeX
    \@ifpackagelater{grffile}{2019/11/01}
    {
      % Do nothing on new versions
    }
    {
      \def\Gread@@xetex#1{%
        \IfFileExists{"\Gin@base".bb}%
        {\Gread@eps{\Gin@base.bb}}%
        {\Gread@@xetex@aux#1}%
      }
    }
    \makeatother
    \usepackage[Export]{adjustbox} % Used to constrain images to a maximum size
    \adjustboxset{max size={0.9\linewidth}{0.9\paperheight}}

    % The hyperref package gives us a pdf with properly built
    % internal navigation ('pdf bookmarks' for the table of contents,
    % internal cross-reference links, web links for URLs, etc.)
    \usepackage{hyperref}
    % The default LaTeX title has an obnoxious amount of whitespace. By default,
    % titling removes some of it. It also provides customization options.
    \usepackage{titling}
    \usepackage{longtable} % longtable support required by pandoc >1.10
    \usepackage{booktabs}  % table support for pandoc > 1.12.2
    \usepackage{array}     % table support for pandoc >= 2.11.3
    \usepackage{calc}      % table minipage width calculation for pandoc >= 2.11.1
    \usepackage[inline]{enumitem} % IRkernel/repr support (it uses the enumerate* environment)
    \usepackage[normalem]{ulem} % ulem is needed to support strikethroughs (\sout)
                                % normalem makes italics be italics, not underlines
    \usepackage{soul}      % strikethrough (\st) support for pandoc >= 3.0.0
    \usepackage{mathrsfs}
    

    
    % Colors for the hyperref package
    \definecolor{urlcolor}{rgb}{0,.145,.698}
    \definecolor{linkcolor}{rgb}{.71,0.21,0.01}
    \definecolor{citecolor}{rgb}{.12,.54,.11}

    % ANSI colors
    \definecolor{ansi-black}{HTML}{3E424D}
    \definecolor{ansi-black-intense}{HTML}{282C36}
    \definecolor{ansi-red}{HTML}{E75C58}
    \definecolor{ansi-red-intense}{HTML}{B22B31}
    \definecolor{ansi-green}{HTML}{00A250}
    \definecolor{ansi-green-intense}{HTML}{007427}
    \definecolor{ansi-yellow}{HTML}{DDB62B}
    \definecolor{ansi-yellow-intense}{HTML}{B27D12}
    \definecolor{ansi-blue}{HTML}{208FFB}
    \definecolor{ansi-blue-intense}{HTML}{0065CA}
    \definecolor{ansi-magenta}{HTML}{D160C4}
    \definecolor{ansi-magenta-intense}{HTML}{A03196}
    \definecolor{ansi-cyan}{HTML}{60C6C8}
    \definecolor{ansi-cyan-intense}{HTML}{258F8F}
    \definecolor{ansi-white}{HTML}{C5C1B4}
    \definecolor{ansi-white-intense}{HTML}{A1A6B2}
    \definecolor{ansi-default-inverse-fg}{HTML}{FFFFFF}
    \definecolor{ansi-default-inverse-bg}{HTML}{000000}

    % common color for the border for error outputs.
    \definecolor{outerrorbackground}{HTML}{FFDFDF}

    % commands and environments needed by pandoc snippets
    % extracted from the output of `pandoc -s`
    \providecommand{\tightlist}{%
      \setlength{\itemsep}{0pt}\setlength{\parskip}{0pt}}
    \DefineVerbatimEnvironment{Highlighting}{Verbatim}{commandchars=\\\{\}}
    % Add ',fontsize=\small' for more characters per line
    \newenvironment{Shaded}{}{}
    \newcommand{\KeywordTok}[1]{\textcolor[rgb]{0.00,0.44,0.13}{\textbf{{#1}}}}
    \newcommand{\DataTypeTok}[1]{\textcolor[rgb]{0.56,0.13,0.00}{{#1}}}
    \newcommand{\DecValTok}[1]{\textcolor[rgb]{0.25,0.63,0.44}{{#1}}}
    \newcommand{\BaseNTok}[1]{\textcolor[rgb]{0.25,0.63,0.44}{{#1}}}
    \newcommand{\FloatTok}[1]{\textcolor[rgb]{0.25,0.63,0.44}{{#1}}}
    \newcommand{\CharTok}[1]{\textcolor[rgb]{0.25,0.44,0.63}{{#1}}}
    \newcommand{\StringTok}[1]{\textcolor[rgb]{0.25,0.44,0.63}{{#1}}}
    \newcommand{\CommentTok}[1]{\textcolor[rgb]{0.38,0.63,0.69}{\textit{{#1}}}}
    \newcommand{\OtherTok}[1]{\textcolor[rgb]{0.00,0.44,0.13}{{#1}}}
    \newcommand{\AlertTok}[1]{\textcolor[rgb]{1.00,0.00,0.00}{\textbf{{#1}}}}
    \newcommand{\FunctionTok}[1]{\textcolor[rgb]{0.02,0.16,0.49}{{#1}}}
    \newcommand{\RegionMarkerTok}[1]{{#1}}
    \newcommand{\ErrorTok}[1]{\textcolor[rgb]{1.00,0.00,0.00}{\textbf{{#1}}}}
    \newcommand{\NormalTok}[1]{{#1}}

    % Additional commands for more recent versions of Pandoc
    \newcommand{\ConstantTok}[1]{\textcolor[rgb]{0.53,0.00,0.00}{{#1}}}
    \newcommand{\SpecialCharTok}[1]{\textcolor[rgb]{0.25,0.44,0.63}{{#1}}}
    \newcommand{\VerbatimStringTok}[1]{\textcolor[rgb]{0.25,0.44,0.63}{{#1}}}
    \newcommand{\SpecialStringTok}[1]{\textcolor[rgb]{0.73,0.40,0.53}{{#1}}}
    \newcommand{\ImportTok}[1]{{#1}}
    \newcommand{\DocumentationTok}[1]{\textcolor[rgb]{0.73,0.13,0.13}{\textit{{#1}}}}
    \newcommand{\AnnotationTok}[1]{\textcolor[rgb]{0.38,0.63,0.69}{\textbf{\textit{{#1}}}}}
    \newcommand{\CommentVarTok}[1]{\textcolor[rgb]{0.38,0.63,0.69}{\textbf{\textit{{#1}}}}}
    \newcommand{\VariableTok}[1]{\textcolor[rgb]{0.10,0.09,0.49}{{#1}}}
    \newcommand{\ControlFlowTok}[1]{\textcolor[rgb]{0.00,0.44,0.13}{\textbf{{#1}}}}
    \newcommand{\OperatorTok}[1]{\textcolor[rgb]{0.40,0.40,0.40}{{#1}}}
    \newcommand{\BuiltInTok}[1]{{#1}}
    \newcommand{\ExtensionTok}[1]{{#1}}
    \newcommand{\PreprocessorTok}[1]{\textcolor[rgb]{0.74,0.48,0.00}{{#1}}}
    \newcommand{\AttributeTok}[1]{\textcolor[rgb]{0.49,0.56,0.16}{{#1}}}
    \newcommand{\InformationTok}[1]{\textcolor[rgb]{0.38,0.63,0.69}{\textbf{\textit{{#1}}}}}
    \newcommand{\WarningTok}[1]{\textcolor[rgb]{0.38,0.63,0.69}{\textbf{\textit{{#1}}}}}
    \makeatletter
    \newsavebox\pandoc@box
    \newcommand*\pandocbounded[1]{%
      \sbox\pandoc@box{#1}%
      % scaling factors for width and height
      \Gscale@div\@tempa\textheight{\dimexpr\ht\pandoc@box+\dp\pandoc@box\relax}%
      \Gscale@div\@tempb\linewidth{\wd\pandoc@box}%
      % select the smaller of both
      \ifdim\@tempb\p@<\@tempa\p@
        \let\@tempa\@tempb
      \fi
      % scaling accordingly (\@tempa < 1)
      \ifdim\@tempa\p@<\p@
        \scalebox{\@tempa}{\usebox\pandoc@box}%
      % scaling not needed, use as it is
      \else
        \usebox{\pandoc@box}%
      \fi
    }
    \makeatother

    % Define a nice break command that doesn't care if a line doesn't already
    % exist.
    \def\br{\hspace*{\fill} \\* }
    % Math Jax compatibility definitions
    \def\gt{>}
    \def\lt{<}
    \let\Oldtex\TeX
    \let\Oldlatex\LaTeX
    \renewcommand{\TeX}{\textrm{\Oldtex}}
    \renewcommand{\LaTeX}{\textrm{\Oldlatex}}
    % Document parameters
    % Document title
    \title{Tratamento\_e\_limpeza\_dados}
    
    
    
    
    
    
    
% Pygments definitions
\makeatletter
\def\PY@reset{\let\PY@it=\relax \let\PY@bf=\relax%
    \let\PY@ul=\relax \let\PY@tc=\relax%
    \let\PY@bc=\relax \let\PY@ff=\relax}
\def\PY@tok#1{\csname PY@tok@#1\endcsname}
\def\PY@toks#1+{\ifx\relax#1\empty\else%
    \PY@tok{#1}\expandafter\PY@toks\fi}
\def\PY@do#1{\PY@bc{\PY@tc{\PY@ul{%
    \PY@it{\PY@bf{\PY@ff{#1}}}}}}}
\def\PY#1#2{\PY@reset\PY@toks#1+\relax+\PY@do{#2}}

\@namedef{PY@tok@w}{\def\PY@tc##1{\textcolor[rgb]{0.73,0.73,0.73}{##1}}}
\@namedef{PY@tok@c}{\let\PY@it=\textit\def\PY@tc##1{\textcolor[rgb]{0.24,0.48,0.48}{##1}}}
\@namedef{PY@tok@cp}{\def\PY@tc##1{\textcolor[rgb]{0.61,0.40,0.00}{##1}}}
\@namedef{PY@tok@k}{\let\PY@bf=\textbf\def\PY@tc##1{\textcolor[rgb]{0.00,0.50,0.00}{##1}}}
\@namedef{PY@tok@kp}{\def\PY@tc##1{\textcolor[rgb]{0.00,0.50,0.00}{##1}}}
\@namedef{PY@tok@kt}{\def\PY@tc##1{\textcolor[rgb]{0.69,0.00,0.25}{##1}}}
\@namedef{PY@tok@o}{\def\PY@tc##1{\textcolor[rgb]{0.40,0.40,0.40}{##1}}}
\@namedef{PY@tok@ow}{\let\PY@bf=\textbf\def\PY@tc##1{\textcolor[rgb]{0.67,0.13,1.00}{##1}}}
\@namedef{PY@tok@nb}{\def\PY@tc##1{\textcolor[rgb]{0.00,0.50,0.00}{##1}}}
\@namedef{PY@tok@nf}{\def\PY@tc##1{\textcolor[rgb]{0.00,0.00,1.00}{##1}}}
\@namedef{PY@tok@nc}{\let\PY@bf=\textbf\def\PY@tc##1{\textcolor[rgb]{0.00,0.00,1.00}{##1}}}
\@namedef{PY@tok@nn}{\let\PY@bf=\textbf\def\PY@tc##1{\textcolor[rgb]{0.00,0.00,1.00}{##1}}}
\@namedef{PY@tok@ne}{\let\PY@bf=\textbf\def\PY@tc##1{\textcolor[rgb]{0.80,0.25,0.22}{##1}}}
\@namedef{PY@tok@nv}{\def\PY@tc##1{\textcolor[rgb]{0.10,0.09,0.49}{##1}}}
\@namedef{PY@tok@no}{\def\PY@tc##1{\textcolor[rgb]{0.53,0.00,0.00}{##1}}}
\@namedef{PY@tok@nl}{\def\PY@tc##1{\textcolor[rgb]{0.46,0.46,0.00}{##1}}}
\@namedef{PY@tok@ni}{\let\PY@bf=\textbf\def\PY@tc##1{\textcolor[rgb]{0.44,0.44,0.44}{##1}}}
\@namedef{PY@tok@na}{\def\PY@tc##1{\textcolor[rgb]{0.41,0.47,0.13}{##1}}}
\@namedef{PY@tok@nt}{\let\PY@bf=\textbf\def\PY@tc##1{\textcolor[rgb]{0.00,0.50,0.00}{##1}}}
\@namedef{PY@tok@nd}{\def\PY@tc##1{\textcolor[rgb]{0.67,0.13,1.00}{##1}}}
\@namedef{PY@tok@s}{\def\PY@tc##1{\textcolor[rgb]{0.73,0.13,0.13}{##1}}}
\@namedef{PY@tok@sd}{\let\PY@it=\textit\def\PY@tc##1{\textcolor[rgb]{0.73,0.13,0.13}{##1}}}
\@namedef{PY@tok@si}{\let\PY@bf=\textbf\def\PY@tc##1{\textcolor[rgb]{0.64,0.35,0.47}{##1}}}
\@namedef{PY@tok@se}{\let\PY@bf=\textbf\def\PY@tc##1{\textcolor[rgb]{0.67,0.36,0.12}{##1}}}
\@namedef{PY@tok@sr}{\def\PY@tc##1{\textcolor[rgb]{0.64,0.35,0.47}{##1}}}
\@namedef{PY@tok@ss}{\def\PY@tc##1{\textcolor[rgb]{0.10,0.09,0.49}{##1}}}
\@namedef{PY@tok@sx}{\def\PY@tc##1{\textcolor[rgb]{0.00,0.50,0.00}{##1}}}
\@namedef{PY@tok@m}{\def\PY@tc##1{\textcolor[rgb]{0.40,0.40,0.40}{##1}}}
\@namedef{PY@tok@gh}{\let\PY@bf=\textbf\def\PY@tc##1{\textcolor[rgb]{0.00,0.00,0.50}{##1}}}
\@namedef{PY@tok@gu}{\let\PY@bf=\textbf\def\PY@tc##1{\textcolor[rgb]{0.50,0.00,0.50}{##1}}}
\@namedef{PY@tok@gd}{\def\PY@tc##1{\textcolor[rgb]{0.63,0.00,0.00}{##1}}}
\@namedef{PY@tok@gi}{\def\PY@tc##1{\textcolor[rgb]{0.00,0.52,0.00}{##1}}}
\@namedef{PY@tok@gr}{\def\PY@tc##1{\textcolor[rgb]{0.89,0.00,0.00}{##1}}}
\@namedef{PY@tok@ge}{\let\PY@it=\textit}
\@namedef{PY@tok@gs}{\let\PY@bf=\textbf}
\@namedef{PY@tok@ges}{\let\PY@bf=\textbf\let\PY@it=\textit}
\@namedef{PY@tok@gp}{\let\PY@bf=\textbf\def\PY@tc##1{\textcolor[rgb]{0.00,0.00,0.50}{##1}}}
\@namedef{PY@tok@go}{\def\PY@tc##1{\textcolor[rgb]{0.44,0.44,0.44}{##1}}}
\@namedef{PY@tok@gt}{\def\PY@tc##1{\textcolor[rgb]{0.00,0.27,0.87}{##1}}}
\@namedef{PY@tok@err}{\def\PY@bc##1{{\setlength{\fboxsep}{\string -\fboxrule}\fcolorbox[rgb]{1.00,0.00,0.00}{1,1,1}{\strut ##1}}}}
\@namedef{PY@tok@kc}{\let\PY@bf=\textbf\def\PY@tc##1{\textcolor[rgb]{0.00,0.50,0.00}{##1}}}
\@namedef{PY@tok@kd}{\let\PY@bf=\textbf\def\PY@tc##1{\textcolor[rgb]{0.00,0.50,0.00}{##1}}}
\@namedef{PY@tok@kn}{\let\PY@bf=\textbf\def\PY@tc##1{\textcolor[rgb]{0.00,0.50,0.00}{##1}}}
\@namedef{PY@tok@kr}{\let\PY@bf=\textbf\def\PY@tc##1{\textcolor[rgb]{0.00,0.50,0.00}{##1}}}
\@namedef{PY@tok@bp}{\def\PY@tc##1{\textcolor[rgb]{0.00,0.50,0.00}{##1}}}
\@namedef{PY@tok@fm}{\def\PY@tc##1{\textcolor[rgb]{0.00,0.00,1.00}{##1}}}
\@namedef{PY@tok@vc}{\def\PY@tc##1{\textcolor[rgb]{0.10,0.09,0.49}{##1}}}
\@namedef{PY@tok@vg}{\def\PY@tc##1{\textcolor[rgb]{0.10,0.09,0.49}{##1}}}
\@namedef{PY@tok@vi}{\def\PY@tc##1{\textcolor[rgb]{0.10,0.09,0.49}{##1}}}
\@namedef{PY@tok@vm}{\def\PY@tc##1{\textcolor[rgb]{0.10,0.09,0.49}{##1}}}
\@namedef{PY@tok@sa}{\def\PY@tc##1{\textcolor[rgb]{0.73,0.13,0.13}{##1}}}
\@namedef{PY@tok@sb}{\def\PY@tc##1{\textcolor[rgb]{0.73,0.13,0.13}{##1}}}
\@namedef{PY@tok@sc}{\def\PY@tc##1{\textcolor[rgb]{0.73,0.13,0.13}{##1}}}
\@namedef{PY@tok@dl}{\def\PY@tc##1{\textcolor[rgb]{0.73,0.13,0.13}{##1}}}
\@namedef{PY@tok@s2}{\def\PY@tc##1{\textcolor[rgb]{0.73,0.13,0.13}{##1}}}
\@namedef{PY@tok@sh}{\def\PY@tc##1{\textcolor[rgb]{0.73,0.13,0.13}{##1}}}
\@namedef{PY@tok@s1}{\def\PY@tc##1{\textcolor[rgb]{0.73,0.13,0.13}{##1}}}
\@namedef{PY@tok@mb}{\def\PY@tc##1{\textcolor[rgb]{0.40,0.40,0.40}{##1}}}
\@namedef{PY@tok@mf}{\def\PY@tc##1{\textcolor[rgb]{0.40,0.40,0.40}{##1}}}
\@namedef{PY@tok@mh}{\def\PY@tc##1{\textcolor[rgb]{0.40,0.40,0.40}{##1}}}
\@namedef{PY@tok@mi}{\def\PY@tc##1{\textcolor[rgb]{0.40,0.40,0.40}{##1}}}
\@namedef{PY@tok@il}{\def\PY@tc##1{\textcolor[rgb]{0.40,0.40,0.40}{##1}}}
\@namedef{PY@tok@mo}{\def\PY@tc##1{\textcolor[rgb]{0.40,0.40,0.40}{##1}}}
\@namedef{PY@tok@ch}{\let\PY@it=\textit\def\PY@tc##1{\textcolor[rgb]{0.24,0.48,0.48}{##1}}}
\@namedef{PY@tok@cm}{\let\PY@it=\textit\def\PY@tc##1{\textcolor[rgb]{0.24,0.48,0.48}{##1}}}
\@namedef{PY@tok@cpf}{\let\PY@it=\textit\def\PY@tc##1{\textcolor[rgb]{0.24,0.48,0.48}{##1}}}
\@namedef{PY@tok@c1}{\let\PY@it=\textit\def\PY@tc##1{\textcolor[rgb]{0.24,0.48,0.48}{##1}}}
\@namedef{PY@tok@cs}{\let\PY@it=\textit\def\PY@tc##1{\textcolor[rgb]{0.24,0.48,0.48}{##1}}}

\def\PYZbs{\char`\\}
\def\PYZus{\char`\_}
\def\PYZob{\char`\{}
\def\PYZcb{\char`\}}
\def\PYZca{\char`\^}
\def\PYZam{\char`\&}
\def\PYZlt{\char`\<}
\def\PYZgt{\char`\>}
\def\PYZsh{\char`\#}
\def\PYZpc{\char`\%}
\def\PYZdl{\char`\$}
\def\PYZhy{\char`\-}
\def\PYZsq{\char`\'}
\def\PYZdq{\char`\"}
\def\PYZti{\char`\~}
% for compatibility with earlier versions
\def\PYZat{@}
\def\PYZlb{[}
\def\PYZrb{]}
\makeatother


    % For linebreaks inside Verbatim environment from package fancyvrb.
    \makeatletter
        \newbox\Wrappedcontinuationbox
        \newbox\Wrappedvisiblespacebox
        \newcommand*\Wrappedvisiblespace {\textcolor{red}{\textvisiblespace}}
        \newcommand*\Wrappedcontinuationsymbol {\textcolor{red}{\llap{\tiny$\m@th\hookrightarrow$}}}
        \newcommand*\Wrappedcontinuationindent {3ex }
        \newcommand*\Wrappedafterbreak {\kern\Wrappedcontinuationindent\copy\Wrappedcontinuationbox}
        % Take advantage of the already applied Pygments mark-up to insert
        % potential linebreaks for TeX processing.
        %        {, <, #, %, $, ' and ": go to next line.
        %        _, }, ^, &, >, - and ~: stay at end of broken line.
        % Use of \textquotesingle for straight quote.
        \newcommand*\Wrappedbreaksatspecials {%
            \def\PYGZus{\discretionary{\char`\_}{\Wrappedafterbreak}{\char`\_}}%
            \def\PYGZob{\discretionary{}{\Wrappedafterbreak\char`\{}{\char`\{}}%
            \def\PYGZcb{\discretionary{\char`\}}{\Wrappedafterbreak}{\char`\}}}%
            \def\PYGZca{\discretionary{\char`\^}{\Wrappedafterbreak}{\char`\^}}%
            \def\PYGZam{\discretionary{\char`\&}{\Wrappedafterbreak}{\char`\&}}%
            \def\PYGZlt{\discretionary{}{\Wrappedafterbreak\char`\<}{\char`\<}}%
            \def\PYGZgt{\discretionary{\char`\>}{\Wrappedafterbreak}{\char`\>}}%
            \def\PYGZsh{\discretionary{}{\Wrappedafterbreak\char`\#}{\char`\#}}%
            \def\PYGZpc{\discretionary{}{\Wrappedafterbreak\char`\%}{\char`\%}}%
            \def\PYGZdl{\discretionary{}{\Wrappedafterbreak\char`\$}{\char`\$}}%
            \def\PYGZhy{\discretionary{\char`\-}{\Wrappedafterbreak}{\char`\-}}%
            \def\PYGZsq{\discretionary{}{\Wrappedafterbreak\textquotesingle}{\textquotesingle}}%
            \def\PYGZdq{\discretionary{}{\Wrappedafterbreak\char`\"}{\char`\"}}%
            \def\PYGZti{\discretionary{\char`\~}{\Wrappedafterbreak}{\char`\~}}%
        }
        % Some characters . , ; ? ! / are not pygmentized.
        % This macro makes them "active" and they will insert potential linebreaks
        \newcommand*\Wrappedbreaksatpunct {%
            \lccode`\~`\.\lowercase{\def~}{\discretionary{\hbox{\char`\.}}{\Wrappedafterbreak}{\hbox{\char`\.}}}%
            \lccode`\~`\,\lowercase{\def~}{\discretionary{\hbox{\char`\,}}{\Wrappedafterbreak}{\hbox{\char`\,}}}%
            \lccode`\~`\;\lowercase{\def~}{\discretionary{\hbox{\char`\;}}{\Wrappedafterbreak}{\hbox{\char`\;}}}%
            \lccode`\~`\:\lowercase{\def~}{\discretionary{\hbox{\char`\:}}{\Wrappedafterbreak}{\hbox{\char`\:}}}%
            \lccode`\~`\?\lowercase{\def~}{\discretionary{\hbox{\char`\?}}{\Wrappedafterbreak}{\hbox{\char`\?}}}%
            \lccode`\~`\!\lowercase{\def~}{\discretionary{\hbox{\char`\!}}{\Wrappedafterbreak}{\hbox{\char`\!}}}%
            \lccode`\~`\/\lowercase{\def~}{\discretionary{\hbox{\char`\/}}{\Wrappedafterbreak}{\hbox{\char`\/}}}%
            \catcode`\.\active
            \catcode`\,\active
            \catcode`\;\active
            \catcode`\:\active
            \catcode`\?\active
            \catcode`\!\active
            \catcode`\/\active
            \lccode`\~`\~
        }
    \makeatother

    \let\OriginalVerbatim=\Verbatim
    \makeatletter
    \renewcommand{\Verbatim}[1][1]{%
        %\parskip\z@skip
        \sbox\Wrappedcontinuationbox {\Wrappedcontinuationsymbol}%
        \sbox\Wrappedvisiblespacebox {\FV@SetupFont\Wrappedvisiblespace}%
        \def\FancyVerbFormatLine ##1{\hsize\linewidth
            \vtop{\raggedright\hyphenpenalty\z@\exhyphenpenalty\z@
                \doublehyphendemerits\z@\finalhyphendemerits\z@
                \strut ##1\strut}%
        }%
        % If the linebreak is at a space, the latter will be displayed as visible
        % space at end of first line, and a continuation symbol starts next line.
        % Stretch/shrink are however usually zero for typewriter font.
        \def\FV@Space {%
            \nobreak\hskip\z@ plus\fontdimen3\font minus\fontdimen4\font
            \discretionary{\copy\Wrappedvisiblespacebox}{\Wrappedafterbreak}
            {\kern\fontdimen2\font}%
        }%

        % Allow breaks at special characters using \PYG... macros.
        \Wrappedbreaksatspecials
        % Breaks at punctuation characters . , ; ? ! and / need catcode=\active
        \OriginalVerbatim[#1,codes*=\Wrappedbreaksatpunct]%
    }
    \makeatother

    % Exact colors from NB
    \definecolor{incolor}{HTML}{303F9F}
    \definecolor{outcolor}{HTML}{D84315}
    \definecolor{cellborder}{HTML}{CFCFCF}
    \definecolor{cellbackground}{HTML}{F7F7F7}

    % prompt
    \makeatletter
    \newcommand{\boxspacing}{\kern\kvtcb@left@rule\kern\kvtcb@boxsep}
    \makeatother
    \newcommand{\prompt}[4]{
        {\ttfamily\llap{{\color{#2}[#3]:\hspace{3pt}#4}}\vspace{-\baselineskip}}
    }
    

    
    % Prevent overflowing lines due to hard-to-break entities
    \sloppy
    % Setup hyperref package
    \hypersetup{
      breaklinks=true,  % so long urls are correctly broken across lines
      colorlinks=true,
      urlcolor=urlcolor,
      linkcolor=linkcolor,
      citecolor=citecolor,
      }
    % Slightly bigger margins than the latex defaults
    
    \geometry{verbose,tmargin=1in,bmargin=1in,lmargin=1in,rmargin=1in}
    
    

\begin{document}
    
    \maketitle
    
    

    
    \subsection{MLOps of Salary Prediction
Classification}\label{mlops-of-salary-prediction-classification}

    \begin{tcolorbox}[breakable, size=fbox, boxrule=1pt, pad at break*=1mm,colback=cellbackground, colframe=cellborder]
\prompt{In}{incolor}{12}{\boxspacing}
\begin{Verbatim}[commandchars=\\\{\}]
\PY{o}{!}pip\PY{+w}{ }install\PY{+w}{ }\PYZhy{}r\PY{+w}{ }requirements.txt
\end{Verbatim}
\end{tcolorbox}

    \begin{Verbatim}[commandchars=\\\{\}]
\^{}C
    \end{Verbatim}

    \subsection{Visualização inicial do
Dataset}\label{visualizauxe7uxe3o-inicial-do-dataset}

Nesta etapa, será realizada uma análise inicial do dataset, incluindo
tamanho, visualização das primeiras linhas, tipos de dados e
estatísticas descritivas

    \begin{tcolorbox}[breakable, size=fbox, boxrule=1pt, pad at break*=1mm,colback=cellbackground, colframe=cellborder]
\prompt{In}{incolor}{9}{\boxspacing}
\begin{Verbatim}[commandchars=\\\{\}]
\PY{k+kn}{import}\PY{+w}{ }\PY{n+nn}{pandas}\PY{+w}{ }\PY{k}{as}\PY{+w}{ }\PY{n+nn}{pd}
\PY{k+kn}{import}\PY{+w}{ }\PY{n+nn}{matplotlib}\PY{n+nn}{.}\PY{n+nn}{pyplot}\PY{+w}{ }\PY{k}{as}\PY{+w}{ }\PY{n+nn}{plt}
\PY{n}{df} \PY{o}{=} \PY{n}{pd}\PY{o}{.}\PY{n}{read\PYZus{}csv}\PY{p}{(}\PY{l+s+s2}{\PYZdq{}}\PY{l+s+s2}{C:/Users/josem/OneDrive/Documentos/MLOPs\PYZus{}Salary\PYZus{}Prediction/salary.csv}\PY{l+s+s2}{\PYZdq{}}\PY{p}{)}
\PY{l+s+sd}{\PYZsq{}\PYZsq{}\PYZsq{}}
\PY{l+s+sd}{print(\PYZdq{}Primeiras linhas do Dataset\PYZdq{})}
\PY{l+s+sd}{print(df.head()) \PYZsh{} Mostra as Primeiras 5 linhas do Dataset}
\PY{l+s+sd}{\PYZsq{}\PYZsq{}\PYZsq{}}
\PY{n+nb}{print}\PY{p}{(}\PY{l+s+s2}{\PYZdq{}}\PY{l+s+se}{\PYZbs{}n}\PY{l+s+s2}{ Formato do Dataset}\PY{l+s+s2}{\PYZdq{}}\PY{p}{)}
\PY{n+nb}{print}\PY{p}{(}\PY{n}{df}\PY{o}{.}\PY{n}{shape}\PY{p}{)} \PY{c+c1}{\PYZsh{} Tamanho do Dataset}

\PY{n+nb}{print}\PY{p}{(}\PY{l+s+s2}{\PYZdq{}}\PY{l+s+se}{\PYZbs{}n}\PY{l+s+s2}{ Informações}\PY{l+s+s2}{\PYZdq{}}\PY{p}{)}
\PY{n}{df}\PY{o}{.}\PY{n}{info}\PY{p}{(}\PY{p}{)} \PY{c+c1}{\PYZsh{} Informações}

\PY{n+nb}{print}\PY{p}{(}\PY{l+s+s2}{\PYZdq{}}\PY{l+s+se}{\PYZbs{}n}\PY{l+s+s2}{ Estatísticas}\PY{l+s+s2}{\PYZdq{}}\PY{p}{)}
\PY{n+nb}{print}\PY{p}{(}\PY{n}{df}\PY{o}{.}\PY{n}{describe}\PY{p}{(}\PY{p}{)}\PY{p}{)}

\PY{n+nb}{print}\PY{p}{(}\PY{l+s+s2}{\PYZdq{}}\PY{l+s+se}{\PYZbs{}n}\PY{l+s+s2}{ Valores Faltantes:}\PY{l+s+s2}{\PYZdq{}}\PY{p}{)}
\PY{n+nb}{print}\PY{p}{(}\PY{n}{df}\PY{o}{.}\PY{n}{isin}\PY{p}{(}\PY{p}{[}\PY{l+s+s2}{\PYZdq{}}\PY{l+s+s2}{?}\PY{l+s+s2}{\PYZdq{}}\PY{p}{]}\PY{p}{)}\PY{o}{.}\PY{n}{sum}\PY{p}{(}\PY{p}{)}\PY{p}{)} \PY{c+c1}{\PYZsh{} valor que pode não está armazenado}
\PY{n+nb}{print}\PY{p}{(}\PY{n}{df}\PY{o}{.}\PY{n}{isnull}\PY{p}{(}\PY{p}{)}\PY{o}{.}\PY{n}{sum}\PY{p}{(}\PY{p}{)}\PY{p}{)}
\PY{n+nb}{print}\PY{p}{(}\PY{l+s+s2}{\PYZdq{}}\PY{l+s+s2}{Distribuição de variável alvo}\PY{l+s+s2}{\PYZdq{}}\PY{p}{)}
\PY{n+nb}{print}\PY{p}{(}\PY{n}{df}\PY{p}{[}\PY{l+s+s1}{\PYZsq{}}\PY{l+s+s1}{salary}\PY{l+s+s1}{\PYZsq{}}\PY{p}{]}\PY{o}{.}\PY{n}{value\PYZus{}counts}\PY{p}{(}\PY{n}{normalize}\PY{o}{=}\PY{k+kc}{True}\PY{p}{)}\PY{p}{)}

\PY{n}{df}\PY{o}{.}\PY{n}{hist}\PY{p}{(}\PY{n}{figsize}\PY{o}{=}\PY{p}{(}\PY{l+m+mi}{12}\PY{p}{,}\PY{l+m+mi}{7}\PY{p}{)}\PY{p}{,} \PY{n}{color}\PY{o}{=}\PY{l+s+s2}{\PYZdq{}}\PY{l+s+s2}{green}\PY{l+s+s2}{\PYZdq{}}\PY{p}{)} 
\PY{n}{plt}\PY{o}{.}\PY{n}{savefig}\PY{p}{(}\PY{l+s+s2}{\PYZdq{}}\PY{l+s+s2}{histogramas.png}\PY{l+s+s2}{\PYZdq{}}\PY{p}{,} \PY{n}{dpi}\PY{o}{=}\PY{l+m+mi}{300}\PY{p}{,} \PY{n}{bbox\PYZus{}inches}\PY{o}{=}\PY{l+s+s1}{\PYZsq{}}\PY{l+s+s1}{tight}\PY{l+s+s1}{\PYZsq{}}\PY{p}{)}
\PY{n}{plt}\PY{o}{.}\PY{n}{show}\PY{p}{(}\PY{p}{)}
\end{Verbatim}
\end{tcolorbox}

    \begin{Verbatim}[commandchars=\\\{\}]

 Formato do Dataset
(32561, 15)

 Informações
<class 'pandas.DataFrame'>
RangeIndex: 32561 entries, 0 to 32560
Data columns (total 15 columns):
 \#   Column          Non-Null Count  Dtype
---  ------          --------------  -----
 0   age             32561 non-null  int64
 1   workclass       32561 non-null  str
 2   fnlwgt          32561 non-null  int64
 3   education       32561 non-null  str
 4   education-num   32561 non-null  int64
 5   marital-status  32561 non-null  str
 6   occupation      32561 non-null  str
 7   relationship    32561 non-null  str
 8   race            32561 non-null  str
 9   sex             32561 non-null  str
 10  capital-gain    32561 non-null  int64
 11  capital-loss    32561 non-null  int64
 12  hours-per-week  32561 non-null  int64
 13  native-country  32561 non-null  str
 14  salary          32561 non-null  str
dtypes: int64(6), str(9)
memory usage: 3.7 MB

 Estatísticas
                age        fnlwgt  education-num  capital-gain  capital-loss  \textbackslash{}
count  32561.000000  3.256100e+04   32561.000000  32561.000000  32561.000000
mean      38.581647  1.897784e+05      10.080679   1077.648844     87.303830
std       13.640433  1.055500e+05       2.572720   7385.292085    402.960219
min       17.000000  1.228500e+04       1.000000      0.000000      0.000000
25\%       28.000000  1.178270e+05       9.000000      0.000000      0.000000
50\%       37.000000  1.783560e+05      10.000000      0.000000      0.000000
75\%       48.000000  2.370510e+05      12.000000      0.000000      0.000000
max       90.000000  1.484705e+06      16.000000  99999.000000   4356.000000

       hours-per-week
count    32561.000000
mean        40.437456
std         12.347429
min          1.000000
25\%         40.000000
50\%         40.000000
75\%         45.000000
max         99.000000

 Valores Faltantes:
age               0
workclass         0
fnlwgt            0
education         0
education-num     0
marital-status    0
occupation        0
relationship      0
race              0
sex               0
capital-gain      0
capital-loss      0
hours-per-week    0
native-country    0
salary            0
dtype: int64
age               0
workclass         0
fnlwgt            0
education         0
education-num     0
marital-status    0
occupation        0
relationship      0
race              0
sex               0
capital-gain      0
capital-loss      0
hours-per-week    0
native-country    0
salary            0
dtype: int64
Distribuição de variável alvo
salary
<=50K    0.75919
>50K     0.24081
Name: proportion, dtype: float64
    \end{Verbatim}

    \begin{center}
    \adjustimage{max size={0.9\linewidth}{0.9\paperheight}}{Tratamento_e_limpeza_dados_files/Tratamento_e_limpeza_dados_3_1.png}
    \end{center}
    { \hspace*{\fill} \\}
    
    \subsection{Análise de Correlação entre
Variáveis}\label{anuxe1lise-de-correlauxe7uxe3o-entre-variuxe1veis}

Nesta etapa, foi realizada a análise de correlação entre todas as
variáveis da base de dados numéricas dentro do dataset de Predição
Salarial, com o objetivo de identificar relações lineares entre
variáveis, possíveis redundâncias de dados e atributos relevantes para o
modelo a fim de fazer o treinamento ser realmente fluido.

O HeatMap abaixo apresenta os coeficientes de correlação de Pearson,
variando entre -1 e 1, onde:

\begin{itemize}
\tightlist
\item
  Valores próximos de \textbf{1} indicam forte correlação positiva\\
\item
  Valores próximos de \textbf{-1} indicam forte correlação negativa\\
\item
  Valores próximos de \textbf{0} indicam baixa ou nenhuma correlação
\end{itemize}

    \begin{tcolorbox}[breakable, size=fbox, boxrule=1pt, pad at break*=1mm,colback=cellbackground, colframe=cellborder]
\prompt{In}{incolor}{11}{\boxspacing}
\begin{Verbatim}[commandchars=\\\{\}]
\PY{k+kn}{import}\PY{+w}{ }\PY{n+nn}{seaborn}\PY{+w}{ }\PY{k}{as}\PY{+w}{ }\PY{n+nn}{sns} \PY{c+c1}{\PYZsh{}importando biblioteca gráfica}
\PY{k+kn}{import}\PY{+w}{ }\PY{n+nn}{matplotlib}\PY{n+nn}{.}\PY{n+nn}{pyplot}\PY{+w}{ }\PY{k}{as}\PY{+w}{ }\PY{n+nn}{plt}

\PY{n}{sns}\PY{o}{.}\PY{n}{heatmap}\PY{p}{(}\PY{n}{df}\PY{o}{.}\PY{n}{select\PYZus{}dtypes}\PY{p}{(}\PY{n}{include}\PY{o}{=}\PY{p}{[}\PY{l+s+s1}{\PYZsq{}}\PY{l+s+s1}{number}\PY{l+s+s1}{\PYZsq{}}\PY{p}{]}\PY{p}{)}\PY{o}{.}\PY{n}{corr}\PY{p}{(}\PY{p}{)}\PY{p}{,} \PY{n}{cmap}\PY{o}{=}\PY{l+s+s1}{\PYZsq{}}\PY{l+s+s1}{Greens}\PY{l+s+s1}{\PYZsq{}}\PY{p}{,} \PY{n}{annot} \PY{o}{=} \PY{k+kc}{True}\PY{p}{)}

\PY{c+c1}{\PYZsh{} depois do gráfico}
\PY{n}{plt}\PY{o}{.}\PY{n}{suptitle}\PY{p}{(}\PY{l+s+s2}{\PYZdq{}}\PY{l+s+s2}{Distribuição das Variáveis Numéricas}\PY{l+s+s2}{\PYZdq{}}\PY{p}{,} \PY{n}{fontsize}\PY{o}{=}\PY{l+m+mi}{16}\PY{p}{)}
\PY{n}{plt}\PY{o}{.}\PY{n}{savefig}\PY{p}{(}\PY{l+s+s2}{\PYZdq{}}\PY{l+s+s2}{grafico\PYZus{}heatmap.png}\PY{l+s+s2}{\PYZdq{}}\PY{p}{,} \PY{n}{dpi}\PY{o}{=}\PY{l+m+mi}{300}\PY{p}{,} \PY{n}{bbox\PYZus{}inches}\PY{o}{=}\PY{l+s+s1}{\PYZsq{}}\PY{l+s+s1}{tight}\PY{l+s+s1}{\PYZsq{}}\PY{p}{)}
\PY{n}{plt}\PY{o}{.}\PY{n}{show}\PY{p}{(}\PY{p}{)}
\end{Verbatim}
\end{tcolorbox}

    \begin{center}
    \adjustimage{max size={0.9\linewidth}{0.9\paperheight}}{Tratamento_e_limpeza_dados_files/Tratamento_e_limpeza_dados_5_0.png}
    \end{center}
    { \hspace*{\fill} \\}
    
    \subsection{Pré-processamento dos
Dados}\label{pruxe9-processamento-dos-dados}

Nesta etapa, será realizado o pré-processamento do dataset com o
objetivo de garantir a qualidade dos dados e melhorar o desempenho do
modelo de classificação. As seguintes técnicas serão aplicadas:

\subsubsection{Remoção de Duplicatas}\label{remouxe7uxe3o-de-duplicatas}

Serão identificados e removidos registros duplicados presentes no
dataset, a fim de evitar redundâncias que possam enviesar o treinamento
do modelo.

\subsubsection{Tratamento de Valores
Faltantes}\label{tratamento-de-valores-faltantes}

Inicialmente, será realizada a análise da proporção de valores ausentes
em cada variável. Colunas que apresentarem mais de 50\% de valores
faltantes serão removidas, por comprometerem a qualidade da informação.
Para as demais, será aplicada imputação de valores, utilizando
estratégias como média, mediana, moda ou valor constante, de acordo com
a distribuição dos dados.

\subsubsection{Tratamento de Outliers}\label{tratamento-de-outliers}

Serão identificados e tratados valores discrepantes (outliers),
utilizando métodos estatísticos como o Intervalo Interquartil (IQR) ou o
Z-score. Essa etapa visa reduzir o impacto de valores extremos que podem
prejudicar a convergência da rede neural.

\subsubsection{Normalização e Padronização dos
Dados}\label{normalizauxe7uxe3o-e-padronizauxe7uxe3o-dos-dados}

As variáveis numéricas serão normalizadas ou padronizadas utilizando
técnicas como Min-Max Scaling ou StandardScaler. Essa etapa é essencial,
pois redes neurais são sensíveis à escala dos dados, e a padronização
contribui para uma convergência mais eficiente durante o treinamento.

    \begin{tcolorbox}[breakable, size=fbox, boxrule=1pt, pad at break*=1mm,colback=cellbackground, colframe=cellborder]
\prompt{In}{incolor}{8}{\boxspacing}
\begin{Verbatim}[commandchars=\\\{\}]
\PY{k+kn}{import}\PY{+w}{ }\PY{n+nn}{pandas}\PY{+w}{ }\PY{k}{as}\PY{+w}{ }\PY{n+nn}{pd}
\PY{k+kn}{import}\PY{+w}{ }\PY{n+nn}{numpy}\PY{+w}{ }\PY{k}{as}\PY{+w}{ }\PY{n+nn}{np}
\PY{k+kn}{from}\PY{+w}{ }\PY{n+nn}{sklearn}\PY{n+nn}{.}\PY{n+nn}{preprocessing}\PY{+w}{ }\PY{k+kn}{import} \PY{n}{StandardScaler}\PY{p}{,} \PY{n}{MinMaxScaler}

\PY{n}{df} \PY{o}{=} \PY{n}{pd}\PY{o}{.}\PY{n}{read\PYZus{}csv}\PY{p}{(}\PY{l+s+s2}{\PYZdq{}}\PY{l+s+s2}{C:/Users/josem/OneDrive/Documentos/MLOPs\PYZus{}Salary\PYZus{}Prediction/salary.csv}\PY{l+s+s2}{\PYZdq{}}\PY{p}{)} \PY{c+c1}{\PYZsh{}Carregamento inicial dos dados}

\PY{n+nb}{print}\PY{p}{(}\PY{l+s+s2}{\PYZdq{}}\PY{l+s+s2}{Shape original:}\PY{l+s+s2}{\PYZdq{}}\PY{p}{,} \PY{n}{df}\PY{o}{.}\PY{n}{shape}\PY{p}{)} \PY{c+c1}{\PYZsh{}Formato antes de remover as duplicatas}

\PY{n}{df} \PY{o}{=} \PY{n}{df}\PY{o}{.}\PY{n}{drop\PYZus{}duplicates}\PY{p}{(}\PY{p}{)} \PY{c+c1}{\PYZsh{} Removendo as duplicatas}
\PY{n+nb}{print}\PY{p}{(}\PY{l+s+s2}{\PYZdq{}}\PY{l+s+s2}{Após remover duplicatas:}\PY{l+s+s2}{\PYZdq{}}\PY{p}{,} \PY{n}{df}\PY{o}{.}\PY{n}{shape}\PY{p}{)} \PY{c+c1}{\PYZsh{}Formato após a remoção das duplicatas}


\PY{c+c1}{\PYZsh{} Substituir \PYZdq{}?\PYZdq{} por NaN (muito comum em datasets)}
\PY{n}{df}\PY{o}{.}\PY{n}{replace}\PY{p}{(}\PY{l+s+s2}{\PYZdq{}}\PY{l+s+s2}{?}\PY{l+s+s2}{\PYZdq{}}\PY{p}{,} \PY{n}{np}\PY{o}{.}\PY{n}{nan}\PY{p}{,} \PY{n}{inplace}\PY{o}{=}\PY{k+kc}{True}\PY{p}{)} \PY{c+c1}{\PYZsh{} Para valores não cadastrados devidamente na base de dados}

\PY{c+c1}{\PYZsh{} Percentual de valores faltantes}
\PY{n}{missing\PYZus{}percent} \PY{o}{=} \PY{n}{df}\PY{o}{.}\PY{n}{isnull}\PY{p}{(}\PY{p}{)}\PY{o}{.}\PY{n}{mean}\PY{p}{(}\PY{p}{)} \PY{o}{*} \PY{l+m+mi}{100}
\PY{n+nb}{print}\PY{p}{(}\PY{l+s+s2}{\PYZdq{}}\PY{l+s+se}{\PYZbs{}n}\PY{l+s+s2}{Percentual de valores faltantes:}\PY{l+s+se}{\PYZbs{}n}\PY{l+s+s2}{\PYZdq{}}\PY{p}{,} \PY{n}{missing\PYZus{}percent}\PY{p}{)}

\PY{c+c1}{\PYZsh{} Remover colunas com mais de 50\PYZpc{} de missing}
\PY{n}{cols\PYZus{}to\PYZus{}drop} \PY{o}{=} \PY{n}{missing\PYZus{}percent}\PY{p}{[}\PY{n}{missing\PYZus{}percent} \PY{o}{\PYZgt{}} \PY{l+m+mi}{50}\PY{p}{]}\PY{o}{.}\PY{n}{index}
\PY{n}{df}\PY{o}{.}\PY{n}{drop}\PY{p}{(}\PY{n}{columns}\PY{o}{=}\PY{n}{cols\PYZus{}to\PYZus{}drop}\PY{p}{,} \PY{n}{inplace}\PY{o}{=}\PY{k+kc}{True}\PY{p}{)}
\PY{n+nb}{print}\PY{p}{(}\PY{l+s+s2}{\PYZdq{}}\PY{l+s+se}{\PYZbs{}n}\PY{l+s+s2}{Colunas removidas (\PYZgt{}50}\PY{l+s+s2}{\PYZpc{}}\PY{l+s+s2}{ missing):}\PY{l+s+s2}{\PYZdq{}}\PY{p}{,} \PY{n+nb}{list}\PY{p}{(}\PY{n}{cols\PYZus{}to\PYZus{}drop}\PY{p}{)}\PY{p}{)}

\PY{c+c1}{\PYZsh{} Separar colunas numéricas e categóricas}
\PY{n}{num\PYZus{}cols} \PY{o}{=} \PY{n}{df}\PY{o}{.}\PY{n}{select\PYZus{}dtypes}\PY{p}{(}\PY{n}{include}\PY{o}{=}\PY{p}{[}\PY{l+s+s1}{\PYZsq{}}\PY{l+s+s1}{int64}\PY{l+s+s1}{\PYZsq{}}\PY{p}{,} \PY{l+s+s1}{\PYZsq{}}\PY{l+s+s1}{float64}\PY{l+s+s1}{\PYZsq{}}\PY{p}{]}\PY{p}{)}\PY{o}{.}\PY{n}{columns}
\PY{n}{cat\PYZus{}cols} \PY{o}{=} \PY{n}{df}\PY{o}{.}\PY{n}{select\PYZus{}dtypes}\PY{p}{(}\PY{n}{include}\PY{o}{=}\PY{p}{[}\PY{l+s+s1}{\PYZsq{}}\PY{l+s+s1}{object}\PY{l+s+s1}{\PYZsq{}}\PY{p}{,} \PY{l+s+s1}{\PYZsq{}}\PY{l+s+s1}{string}\PY{l+s+s1}{\PYZsq{}}\PY{p}{]}\PY{p}{)}\PY{o}{.}\PY{n}{columns}

\PY{c+c1}{\PYZsh{} Imputação numérica (mediana)}
\PY{k}{for} \PY{n}{col} \PY{o+ow}{in} \PY{n}{num\PYZus{}cols}\PY{p}{:}
    \PY{n}{df}\PY{p}{[}\PY{n}{col}\PY{p}{]} \PY{o}{=} \PY{n}{df}\PY{p}{[}\PY{n}{col}\PY{p}{]}\PY{o}{.}\PY{n}{fillna}\PY{p}{(}\PY{n}{df}\PY{p}{[}\PY{n}{col}\PY{p}{]}\PY{o}{.}\PY{n}{median}\PY{p}{(}\PY{p}{)}\PY{p}{)}

\PY{c+c1}{\PYZsh{} Imputação categórica (moda)}
\PY{k}{for} \PY{n}{col} \PY{o+ow}{in} \PY{n}{cat\PYZus{}cols}\PY{p}{:}
    \PY{n}{df}\PY{p}{[}\PY{n}{col}\PY{p}{]} \PY{o}{=} \PY{n}{df}\PY{p}{[}\PY{n}{col}\PY{p}{]}\PY{o}{.}\PY{n}{fillna}\PY{p}{(}\PY{n}{df}\PY{p}{[}\PY{n}{col}\PY{p}{]}\PY{o}{.}\PY{n}{mode}\PY{p}{(}\PY{p}{)}\PY{p}{[}\PY{l+m+mi}{0}\PY{p}{]}\PY{p}{)}

\PY{n+nb}{print}\PY{p}{(}\PY{l+s+s2}{\PYZdq{}}\PY{l+s+se}{\PYZbs{}n}\PY{l+s+s2}{Valores faltantes após tratamento:}\PY{l+s+se}{\PYZbs{}n}\PY{l+s+s2}{\PYZdq{}}\PY{p}{,} \PY{n}{df}\PY{o}{.}\PY{n}{isnull}\PY{p}{(}\PY{p}{)}\PY{o}{.}\PY{n}{sum}\PY{p}{(}\PY{p}{)}\PY{p}{)}

\PY{c+c1}{\PYZsh{} =========================}
\PY{c+c1}{\PYZsh{} 4. Tratamento de outliers (IQR)}
\PY{c+c1}{\PYZsh{} =========================}

\PY{k}{def}\PY{+w}{ }\PY{n+nf}{remove\PYZus{}outliers\PYZus{}iqr}\PY{p}{(}\PY{n}{df}\PY{p}{,} \PY{n}{cols}\PY{p}{)}\PY{p}{:}
    \PY{n}{df\PYZus{}clean} \PY{o}{=} \PY{n}{df}\PY{o}{.}\PY{n}{copy}\PY{p}{(}\PY{p}{)}
    \PY{k}{for} \PY{n}{col} \PY{o+ow}{in} \PY{n}{cols}\PY{p}{:}
        \PY{n}{Q1} \PY{o}{=} \PY{n}{df\PYZus{}clean}\PY{p}{[}\PY{n}{col}\PY{p}{]}\PY{o}{.}\PY{n}{quantile}\PY{p}{(}\PY{l+m+mf}{0.25}\PY{p}{)}
        \PY{n}{Q3} \PY{o}{=} \PY{n}{df\PYZus{}clean}\PY{p}{[}\PY{n}{col}\PY{p}{]}\PY{o}{.}\PY{n}{quantile}\PY{p}{(}\PY{l+m+mf}{0.75}\PY{p}{)}
        \PY{n}{IQR} \PY{o}{=} \PY{n}{Q3} \PY{o}{\PYZhy{}} \PY{n}{Q1}

        \PY{n}{lower} \PY{o}{=} \PY{n}{Q1} \PY{o}{\PYZhy{}} \PY{l+m+mf}{1.5} \PY{o}{*} \PY{n}{IQR}
        \PY{n}{upper} \PY{o}{=} \PY{n}{Q3} \PY{o}{+} \PY{l+m+mf}{1.5} \PY{o}{*} \PY{n}{IQR}

        \PY{n}{df\PYZus{}clean} \PY{o}{=} \PY{n}{df\PYZus{}clean}\PY{p}{[}\PY{p}{(}\PY{n}{df\PYZus{}clean}\PY{p}{[}\PY{n}{col}\PY{p}{]} \PY{o}{\PYZgt{}}\PY{o}{=} \PY{n}{lower}\PY{p}{)} \PY{o}{\PYZam{}} \PY{p}{(}\PY{n}{df\PYZus{}clean}\PY{p}{[}\PY{n}{col}\PY{p}{]} \PY{o}{\PYZlt{}}\PY{o}{=} \PY{n}{upper}\PY{p}{)}\PY{p}{]}

    \PY{k}{return} \PY{n}{df\PYZus{}clean}

\PY{n}{df} \PY{o}{=} \PY{n}{remove\PYZus{}outliers\PYZus{}iqr}\PY{p}{(}\PY{n}{df}\PY{p}{,} \PY{n}{num\PYZus{}cols}\PY{p}{)}
\PY{n+nb}{print}\PY{p}{(}\PY{l+s+s2}{\PYZdq{}}\PY{l+s+s2}{Após remoção de outliers:}\PY{l+s+s2}{\PYZdq{}}\PY{p}{,} \PY{n}{df}\PY{o}{.}\PY{n}{shape}\PY{p}{)}

\PY{c+c1}{\PYZsh{} =========================}
\PY{c+c1}{\PYZsh{} 5. Normalização / Padronização}
\PY{c+c1}{\PYZsh{} =========================}

\PY{c+c1}{\PYZsh{} Escolha: StandardScaler (recomendado para MLP)}
\PY{n}{scaler} \PY{o}{=} \PY{n}{StandardScaler}\PY{p}{(}\PY{p}{)}

\PY{n}{df}\PY{p}{[}\PY{n}{num\PYZus{}cols}\PY{p}{]} \PY{o}{=} \PY{n}{scaler}\PY{o}{.}\PY{n}{fit\PYZus{}transform}\PY{p}{(}\PY{n}{df}\PY{p}{[}\PY{n}{num\PYZus{}cols}\PY{p}{]}\PY{p}{)}

\PY{n+nb}{print}\PY{p}{(}\PY{l+s+s2}{\PYZdq{}}\PY{l+s+se}{\PYZbs{}n}\PY{l+s+s2}{Dados normalizados com StandardScaler.}\PY{l+s+s2}{\PYZdq{}}\PY{p}{)}

\PY{c+c1}{\PYZsh{} =========================}
\PY{c+c1}{\PYZsh{} Resultado final}
\PY{c+c1}{\PYZsh{} =========================}
\PY{n+nb}{print}\PY{p}{(}\PY{l+s+s2}{\PYZdq{}}\PY{l+s+se}{\PYZbs{}n}\PY{l+s+s2}{Shape final:}\PY{l+s+s2}{\PYZdq{}}\PY{p}{,} \PY{n}{df}\PY{o}{.}\PY{n}{shape}\PY{p}{)}
\PY{n+nb}{print}\PY{p}{(}\PY{l+s+s2}{\PYZdq{}}\PY{l+s+se}{\PYZbs{}n}\PY{l+s+s2}{Preview final:}\PY{l+s+s2}{\PYZdq{}}\PY{p}{)}
\PY{n+nb}{print}\PY{p}{(}\PY{n}{df}\PY{o}{.}\PY{n}{head}\PY{p}{(}\PY{p}{)}\PY{p}{)}
\end{Verbatim}
\end{tcolorbox}

    \begin{Verbatim}[commandchars=\\\{\}]
Shape original: (32561, 15)
Após remover duplicatas: (32537, 15)

Percentual de valores faltantes:
 age               0.0
workclass         0.0
fnlwgt            0.0
education         0.0
education-num     0.0
marital-status    0.0
occupation        0.0
relationship      0.0
race              0.0
sex               0.0
capital-gain      0.0
capital-loss      0.0
hours-per-week    0.0
native-country    0.0
salary            0.0
dtype: float64

Colunas removidas (>50\% missing): []

Valores faltantes após tratamento:
 age               0
workclass         0
fnlwgt            0
education         0
education-num     0
marital-status    0
occupation        0
relationship      0
race              0
sex               0
capital-gain      0
capital-loss      0
hours-per-week    0
native-country    0
salary            0
dtype: int64
Após remoção de outliers: (18991, 15)

Dados normalizados com StandardScaler.

Shape final: (18991, 15)

Preview final:
        age          workclass    fnlwgt   education  education-num  \textbackslash{}
2 -0.004199            Private  0.401116     HS-grad      -0.579081
3  1.243037            Private  0.620159        11th      -1.508342
4 -0.835689            Private  1.810833   Bachelors       1.279439
5 -0.087348            Private  1.192725     Masters       1.744070
7  1.159888   Self-emp-not-inc  0.332170     HS-grad      -0.579081

        marital-status          occupation    relationship    race      sex  \textbackslash{}
2             Divorced   Handlers-cleaners   Not-in-family   White     Male
3   Married-civ-spouse   Handlers-cleaners         Husband   Black     Male
4   Married-civ-spouse      Prof-specialty            Wife   Black   Female
5   Married-civ-spouse     Exec-managerial            Wife   White   Female
7   Married-civ-spouse     Exec-managerial         Husband   White     Male

   capital-gain  capital-loss  hours-per-week  native-country  salary
2           0.0           0.0       -0.372235   United-States   <=50K
3           0.0           0.0       -0.372235   United-States   <=50K
4           0.0           0.0       -0.372235            Cuba   <=50K
5           0.0           0.0       -0.372235   United-States   <=50K
7           0.0           0.0        0.893451   United-States    >50K
    \end{Verbatim}

    \subsection{3. Seleção de Variáveis (Feature
Selection)}\label{seleuxe7uxe3o-de-variuxe1veis-feature-selection}

Nesta etapa, aplicamos um processo de filtragem e ranking para garantir
que apenas os atributos mais relevantes e não redundantes sejam
utilizados no treinamento da rede neural MLP. Um conjunto de dados
enxuto e estatisticamente justificado reduz o tempo de treinamento e
evita o overfitting.

\subsection{3.1. Análise de Colinearidade via VIF (Variance Inflation
Factor)}\label{anuxe1lise-de-colinearidade-via-vif-variance-inflation-factor}

O primeiro passo é identificar a multicolinearidade. Quando duas ou mais
variáveis são altamente correlacionadas, elas fornecem informações
redundantes, o que pode causar instabilidade nos pesos da rede
neural.Critério de Exclusão: Variáveis com VIF \textgreater{} 10 são
removidas.Justificativa Técnica: A remoção de variáveis com alto VIF
garante que a matriz de dados seja bem condicionada, facilitando a
convergência do algoritmo de otimização durante o treinamento da
MLP.Nota: Variáveis com VIF ``infinito'' indicam colinearidade perfeita
e foram as primeiras a serem descartadas.

\subsection{3.2. Metodologia de Ranking
Triplo}\label{metodologia-de-ranking-triplo}

Para uma seleção robusta, foram combinadas três perspectivas diferentes
de importância de variáveis:Correlação de Pearson (Estatística Linear):
Mede a força da relação linear entre o atributo e a variável alvo
(salary).Mutual Information (Dependência Não-Linear): Diferente da
correlação, a Informação Mútua captura qualquer tipo de dependência
estatística, sendo essencial para redes neurais que modelam relações
complexas.Random Forest Importance (Baseado em Árvore): Avalia o quanto
cada variável contribui para a redução da impureza (Gini/Entropia) na
tomada de decisão de um conjunto de árvores.3.3. Tabela Comparativa de
ImportânciaAbaixo, os resultados consolidados. Todos os scores foram
normalizados na escala \([0, 1]\) para permitir uma média justa (Score
Final)

    \begin{tcolorbox}[breakable, size=fbox, boxrule=1pt, pad at break*=1mm,colback=cellbackground, colframe=cellborder]
\prompt{In}{incolor}{ }{\boxspacing}
\begin{Verbatim}[commandchars=\\\{\}]
\PY{k+kn}{import}\PY{+w}{ }\PY{n+nn}{pandas}\PY{+w}{ }\PY{k}{as}\PY{+w}{ }\PY{n+nn}{pd}
\PY{k+kn}{import}\PY{+w}{ }\PY{n+nn}{numpy}\PY{+w}{ }\PY{k}{as}\PY{+w}{ }\PY{n+nn}{np}
\PY{k+kn}{from}\PY{+w}{ }\PY{n+nn}{sklearn}\PY{n+nn}{.}\PY{n+nn}{feature\PYZus{}selection}\PY{+w}{ }\PY{k+kn}{import} \PY{n}{mutual\PYZus{}info\PYZus{}classif}
\PY{k+kn}{from}\PY{+w}{ }\PY{n+nn}{sklearn}\PY{n+nn}{.}\PY{n+nn}{ensemble}\PY{+w}{ }\PY{k+kn}{import} \PY{n}{RandomForestClassifier}
\PY{k+kn}{from}\PY{+w}{ }\PY{n+nn}{sklearn}\PY{n+nn}{.}\PY{n+nn}{preprocessing}\PY{+w}{ }\PY{k+kn}{import} \PY{n}{MinMaxScaler}
\PY{k+kn}{from}\PY{+w}{ }\PY{n+nn}{statsmodels}\PY{n+nn}{.}\PY{n+nn}{stats}\PY{n+nn}{.}\PY{n+nn}{outliers\PYZus{}influence}\PY{+w}{ }\PY{k+kn}{import} \PY{n}{variance\PYZus{}inflation\PYZus{}factor}

\PY{c+c1}{\PYZsh{} 1. PREPARAÇÃO E LIMPEZA TOTAL}
\PY{c+c1}{\PYZsh{} \PYZhy{}\PYZhy{}\PYZhy{}\PYZhy{}\PYZhy{}\PYZhy{}\PYZhy{}\PYZhy{}\PYZhy{}\PYZhy{}\PYZhy{}\PYZhy{}\PYZhy{}\PYZhy{}\PYZhy{}\PYZhy{}\PYZhy{}\PYZhy{}\PYZhy{}\PYZhy{}\PYZhy{}\PYZhy{}\PYZhy{}\PYZhy{}\PYZhy{}\PYZhy{}\PYZhy{}\PYZhy{}\PYZhy{}\PYZhy{}\PYZhy{}\PYZhy{}\PYZhy{}\PYZhy{}\PYZhy{}\PYZhy{}\PYZhy{}\PYZhy{}\PYZhy{}\PYZhy{}\PYZhy{}\PYZhy{}\PYZhy{}\PYZhy{}\PYZhy{}\PYZhy{}\PYZhy{}\PYZhy{}\PYZhy{}\PYZhy{}\PYZhy{}\PYZhy{}\PYZhy{}\PYZhy{}\PYZhy{}\PYZhy{}\PYZhy{}}
\PY{c+c1}{\PYZsh{} Criando uma cópia para não afetar o df original}
\PY{n}{X\PYZus{}working} \PY{o}{=} \PY{n}{df}\PY{o}{.}\PY{n}{drop}\PY{p}{(}\PY{n}{columns}\PY{o}{=}\PY{p}{[}\PY{l+s+s1}{\PYZsq{}}\PY{l+s+s1}{salary}\PY{l+s+s1}{\PYZsq{}}\PY{p}{]}\PY{p}{)}\PY{o}{.}\PY{n}{copy}\PY{p}{(}\PY{p}{)}
\PY{n}{y} \PY{o}{=} \PY{n}{df}\PY{p}{[}\PY{l+s+s1}{\PYZsq{}}\PY{l+s+s1}{salary}\PY{l+s+s1}{\PYZsq{}}\PY{p}{]}

\PY{c+c1}{\PYZsh{} Converte variáveis categóricas (object/string) em colunas numéricas (0 e 1)}
\PY{n}{X\PYZus{}encoded} \PY{o}{=} \PY{n}{pd}\PY{o}{.}\PY{n}{get\PYZus{}dummies}\PY{p}{(}\PY{n}{X\PYZus{}working}\PY{p}{,} \PY{n}{drop\PYZus{}first}\PY{o}{=}\PY{k+kc}{True}\PY{p}{)}

\PY{c+c1}{\PYZsh{} FORÇAR TUDO PARA FLOAT: Isso resolve o erro \PYZsq{}ufunc isfinite\PYZsq{}}
\PY{n}{X\PYZus{}encoded} \PY{o}{=} \PY{n}{X\PYZus{}encoded}\PY{o}{.}\PY{n}{astype}\PY{p}{(}\PY{n+nb}{float}\PY{p}{)}

\PY{c+c1}{\PYZsh{} LIMPEZA DE NANS E INFINITOS: O VIF não aceita nenhum valor ausente}
\PY{n}{X\PYZus{}encoded} \PY{o}{=} \PY{n}{X\PYZus{}encoded}\PY{o}{.}\PY{n}{replace}\PY{p}{(}\PY{p}{[}\PY{n}{np}\PY{o}{.}\PY{n}{inf}\PY{p}{,} \PY{o}{\PYZhy{}}\PY{n}{np}\PY{o}{.}\PY{n}{inf}\PY{p}{]}\PY{p}{,} \PY{n}{np}\PY{o}{.}\PY{n}{nan}\PY{p}{)}\PY{o}{.}\PY{n}{dropna}\PY{p}{(}\PY{p}{)}
\PY{c+c1}{\PYZsh{} Ajustar o y para ter o mesmo índice caso linhas tenham sido removidas}
\PY{n}{y} \PY{o}{=} \PY{n}{y}\PY{o}{.}\PY{n}{loc}\PY{p}{[}\PY{n}{X\PYZus{}encoded}\PY{o}{.}\PY{n}{index}\PY{p}{]}

\PY{c+c1}{\PYZsh{} 2. CÁLCULO DO VIF (DIRETRIZ: Justificar exclusão VIF \PYZgt{} 10)}
\PY{c+c1}{\PYZsh{} \PYZhy{}\PYZhy{}\PYZhy{}\PYZhy{}\PYZhy{}\PYZhy{}\PYZhy{}\PYZhy{}\PYZhy{}\PYZhy{}\PYZhy{}\PYZhy{}\PYZhy{}\PYZhy{}\PYZhy{}\PYZhy{}\PYZhy{}\PYZhy{}\PYZhy{}\PYZhy{}\PYZhy{}\PYZhy{}\PYZhy{}\PYZhy{}\PYZhy{}\PYZhy{}\PYZhy{}\PYZhy{}\PYZhy{}\PYZhy{}\PYZhy{}\PYZhy{}\PYZhy{}\PYZhy{}\PYZhy{}\PYZhy{}\PYZhy{}\PYZhy{}\PYZhy{}\PYZhy{}\PYZhy{}\PYZhy{}\PYZhy{}\PYZhy{}\PYZhy{}\PYZhy{}\PYZhy{}\PYZhy{}\PYZhy{}\PYZhy{}\PYZhy{}\PYZhy{}\PYZhy{}\PYZhy{}\PYZhy{}\PYZhy{}\PYZhy{}}
\PY{k}{def}\PY{+w}{ }\PY{n+nf}{calculate\PYZus{}vif}\PY{p}{(}\PY{n}{X\PYZus{}data}\PY{p}{)}\PY{p}{:}
    \PY{c+c1}{\PYZsh{} Adiciona constante para o cálculo estatístico correto do intercepto}
    \PY{n}{X\PYZus{}const} \PY{o}{=} \PY{n}{X\PYZus{}data}\PY{o}{.}\PY{n}{assign}\PY{p}{(}\PY{n}{const}\PY{o}{=}\PY{l+m+mf}{1.0}\PY{p}{)}
    \PY{n}{vif\PYZus{}data} \PY{o}{=} \PY{n}{pd}\PY{o}{.}\PY{n}{DataFrame}\PY{p}{(}\PY{p}{)}
    \PY{n}{vif\PYZus{}data}\PY{p}{[}\PY{l+s+s2}{\PYZdq{}}\PY{l+s+s2}{feature}\PY{l+s+s2}{\PYZdq{}}\PY{p}{]} \PY{o}{=} \PY{n}{X\PYZus{}const}\PY{o}{.}\PY{n}{columns}
    \PY{n}{vif\PYZus{}data}\PY{p}{[}\PY{l+s+s2}{\PYZdq{}}\PY{l+s+s2}{VIF}\PY{l+s+s2}{\PYZdq{}}\PY{p}{]} \PY{o}{=} \PY{p}{[}\PY{n}{variance\PYZus{}inflation\PYZus{}factor}\PY{p}{(}\PY{n}{X\PYZus{}const}\PY{o}{.}\PY{n}{values}\PY{p}{,} \PY{n}{i}\PY{p}{)} \PY{k}{for} \PY{n}{i} \PY{o+ow}{in} \PY{n+nb}{range}\PY{p}{(}\PY{n+nb}{len}\PY{p}{(}\PY{n}{X\PYZus{}const}\PY{o}{.}\PY{n}{columns}\PY{p}{)}\PY{p}{)}\PY{p}{]}
    \PY{k}{return} \PY{n}{vif\PYZus{}data}\PY{p}{[}\PY{n}{vif\PYZus{}data}\PY{p}{[}\PY{l+s+s2}{\PYZdq{}}\PY{l+s+s2}{feature}\PY{l+s+s2}{\PYZdq{}}\PY{p}{]} \PY{o}{!=} \PY{l+s+s2}{\PYZdq{}}\PY{l+s+s2}{const}\PY{l+s+s2}{\PYZdq{}}\PY{p}{]}

\PY{n+nb}{print}\PY{p}{(}\PY{l+s+s2}{\PYZdq{}}\PY{l+s+s2}{Calculando VIF...}\PY{l+s+s2}{\PYZdq{}}\PY{p}{)}
\PY{n}{vif\PYZus{}df} \PY{o}{=} \PY{n}{calculate\PYZus{}vif}\PY{p}{(}\PY{n}{X\PYZus{}encoded}\PY{p}{)}

\PY{c+c1}{\PYZsh{} Identificando features para exclusão conforme as diretrizes }
\PY{n}{high\PYZus{}vif\PYZus{}features} \PY{o}{=} \PY{n}{vif\PYZus{}df}\PY{p}{[}\PY{n}{vif\PYZus{}df}\PY{p}{[}\PY{l+s+s2}{\PYZdq{}}\PY{l+s+s2}{VIF}\PY{l+s+s2}{\PYZdq{}}\PY{p}{]} \PY{o}{\PYZgt{}} \PY{l+m+mi}{10}\PY{p}{]}\PY{p}{[}\PY{l+s+s2}{\PYZdq{}}\PY{l+s+s2}{feature}\PY{l+s+s2}{\PYZdq{}}\PY{p}{]}\PY{o}{.}\PY{n}{tolist}\PY{p}{(}\PY{p}{)}
\PY{n}{X\PYZus{}final} \PY{o}{=} \PY{n}{X\PYZus{}encoded}\PY{o}{.}\PY{n}{drop}\PY{p}{(}\PY{n}{columns}\PY{o}{=}\PY{n}{high\PYZus{}vif\PYZus{}features}\PY{p}{)}

\PY{n+nb}{print}\PY{p}{(}\PY{l+s+sa}{f}\PY{l+s+s2}{\PYZdq{}}\PY{l+s+se}{\PYZbs{}n}\PY{l+s+s2}{Variáveis removidas (VIF \PYZgt{} 10): }\PY{l+s+si}{\PYZob{}}\PY{n}{high\PYZus{}vif\PYZus{}features}\PY{l+s+si}{\PYZcb{}}\PY{l+s+s2}{\PYZdq{}}\PY{p}{)}

\PY{c+c1}{\PYZsh{} 3. RANKING TRI\PYZhy{}MÉTODO (Pearson, Mutual Info, Random Forest)}
\PY{c+c1}{\PYZsh{} \PYZhy{}\PYZhy{}\PYZhy{}\PYZhy{}\PYZhy{}\PYZhy{}\PYZhy{}\PYZhy{}\PYZhy{}\PYZhy{}\PYZhy{}\PYZhy{}\PYZhy{}\PYZhy{}\PYZhy{}\PYZhy{}\PYZhy{}\PYZhy{}\PYZhy{}\PYZhy{}\PYZhy{}\PYZhy{}\PYZhy{}\PYZhy{}\PYZhy{}\PYZhy{}\PYZhy{}\PYZhy{}\PYZhy{}\PYZhy{}\PYZhy{}\PYZhy{}\PYZhy{}\PYZhy{}\PYZhy{}\PYZhy{}\PYZhy{}\PYZhy{}\PYZhy{}\PYZhy{}\PYZhy{}\PYZhy{}\PYZhy{}\PYZhy{}\PYZhy{}\PYZhy{}\PYZhy{}\PYZhy{}\PYZhy{}\PYZhy{}\PYZhy{}\PYZhy{}\PYZhy{}\PYZhy{}\PYZhy{}\PYZhy{}\PYZhy{}}
\PY{n+nb}{print}\PY{p}{(}\PY{l+s+s2}{\PYZdq{}}\PY{l+s+s2}{Gerando ranking triplo...}\PY{l+s+s2}{\PYZdq{}}\PY{p}{)}

\PY{c+c1}{\PYZsh{} Método 1: Pearson (Linear)}
\PY{n}{pearson\PYZus{}scores} \PY{o}{=} \PY{n}{X\PYZus{}final}\PY{o}{.}\PY{n}{corrwith}\PY{p}{(}\PY{n}{y}\PY{p}{)}\PY{o}{.}\PY{n}{abs}\PY{p}{(}\PY{p}{)}\PY{o}{.}\PY{n}{values}

\PY{c+c1}{\PYZsh{} Método 2: Mutual Information (Não\PYZhy{}linear) [cite: 31]}
\PY{n}{mi\PYZus{}scores} \PY{o}{=} \PY{n}{mutual\PYZus{}info\PYZus{}classif}\PY{p}{(}\PY{n}{X\PYZus{}final}\PY{p}{,} \PY{n}{y}\PY{p}{,} \PY{n}{random\PYZus{}state}\PY{o}{=}\PY{l+m+mi}{42}\PY{p}{)}

\PY{c+c1}{\PYZsh{} Método 3: Random Forest (Baseado em Árvore) [cite: 32]}
\PY{n}{rf} \PY{o}{=} \PY{n}{RandomForestClassifier}\PY{p}{(}\PY{n}{n\PYZus{}estimators}\PY{o}{=}\PY{l+m+mi}{100}\PY{p}{,} \PY{n}{random\PYZus{}state}\PY{o}{=}\PY{l+m+mi}{42}\PY{p}{)}
\PY{n}{rf}\PY{o}{.}\PY{n}{fit}\PY{p}{(}\PY{n}{X\PYZus{}final}\PY{p}{,} \PY{n}{y}\PY{p}{)}
\PY{n}{rf\PYZus{}scores} \PY{o}{=} \PY{n}{rf}\PY{o}{.}\PY{n}{feature\PYZus{}importances\PYZus{}}

\PY{c+c1}{\PYZsh{} 4. TABELA COMPARATIVA (EXIGÊNCIA DO RELATÓRIO [cite: 33])}
\PY{c+c1}{\PYZsh{} \PYZhy{}\PYZhy{}\PYZhy{}\PYZhy{}\PYZhy{}\PYZhy{}\PYZhy{}\PYZhy{}\PYZhy{}\PYZhy{}\PYZhy{}\PYZhy{}\PYZhy{}\PYZhy{}\PYZhy{}\PYZhy{}\PYZhy{}\PYZhy{}\PYZhy{}\PYZhy{}\PYZhy{}\PYZhy{}\PYZhy{}\PYZhy{}\PYZhy{}\PYZhy{}\PYZhy{}\PYZhy{}\PYZhy{}\PYZhy{}\PYZhy{}\PYZhy{}\PYZhy{}\PYZhy{}\PYZhy{}\PYZhy{}\PYZhy{}\PYZhy{}\PYZhy{}\PYZhy{}\PYZhy{}\PYZhy{}\PYZhy{}\PYZhy{}\PYZhy{}\PYZhy{}\PYZhy{}\PYZhy{}\PYZhy{}\PYZhy{}\PYZhy{}\PYZhy{}\PYZhy{}\PYZhy{}\PYZhy{}\PYZhy{}\PYZhy{}}
\PY{n}{scaler} \PY{o}{=} \PY{n}{MinMaxScaler}\PY{p}{(}\PY{p}{)}
\PY{n}{rankings} \PY{o}{=} \PY{n}{pd}\PY{o}{.}\PY{n}{DataFrame}\PY{p}{(}\PY{p}{\PYZob{}}
    \PY{l+s+s1}{\PYZsq{}}\PY{l+s+s1}{Feature}\PY{l+s+s1}{\PYZsq{}}\PY{p}{:} \PY{n}{X\PYZus{}final}\PY{o}{.}\PY{n}{columns}\PY{p}{,}
    \PY{l+s+s1}{\PYZsq{}}\PY{l+s+s1}{Pearson}\PY{l+s+s1}{\PYZsq{}}\PY{p}{:} \PY{n}{scaler}\PY{o}{.}\PY{n}{fit\PYZus{}transform}\PY{p}{(}\PY{n}{pearson\PYZus{}scores}\PY{o}{.}\PY{n}{reshape}\PY{p}{(}\PY{o}{\PYZhy{}}\PY{l+m+mi}{1}\PY{p}{,} \PY{l+m+mi}{1}\PY{p}{)}\PY{p}{)}\PY{o}{.}\PY{n}{flatten}\PY{p}{(}\PY{p}{)}\PY{p}{,}
    \PY{l+s+s1}{\PYZsq{}}\PY{l+s+s1}{Mutual\PYZus{}Info}\PY{l+s+s1}{\PYZsq{}}\PY{p}{:} \PY{n}{scaler}\PY{o}{.}\PY{n}{fit\PYZus{}transform}\PY{p}{(}\PY{n}{mi\PYZus{}scores}\PY{o}{.}\PY{n}{reshape}\PY{p}{(}\PY{o}{\PYZhy{}}\PY{l+m+mi}{1}\PY{p}{,} \PY{l+m+mi}{1}\PY{p}{)}\PY{p}{)}\PY{o}{.}\PY{n}{flatten}\PY{p}{(}\PY{p}{)}\PY{p}{,}
    \PY{l+s+s1}{\PYZsq{}}\PY{l+s+s1}{Random\PYZus{}Forest}\PY{l+s+s1}{\PYZsq{}}\PY{p}{:} \PY{n}{scaler}\PY{o}{.}\PY{n}{fit\PYZus{}transform}\PY{p}{(}\PY{n}{rf\PYZus{}scores}\PY{o}{.}\PY{n}{reshape}\PY{p}{(}\PY{o}{\PYZhy{}}\PY{l+m+mi}{1}\PY{p}{,} \PY{l+m+mi}{1}\PY{p}{)}\PY{p}{)}\PY{o}{.}\PY{n}{flatten}\PY{p}{(}\PY{p}{)}
\PY{p}{\PYZcb{}}\PY{p}{)}

\PY{n}{rankings}\PY{p}{[}\PY{l+s+s1}{\PYZsq{}}\PY{l+s+s1}{Score\PYZus{}Final}\PY{l+s+s1}{\PYZsq{}}\PY{p}{]} \PY{o}{=} \PY{n}{rankings}\PY{p}{[}\PY{p}{[}\PY{l+s+s1}{\PYZsq{}}\PY{l+s+s1}{Pearson}\PY{l+s+s1}{\PYZsq{}}\PY{p}{,} \PY{l+s+s1}{\PYZsq{}}\PY{l+s+s1}{Mutual\PYZus{}Info}\PY{l+s+s1}{\PYZsq{}}\PY{p}{,} \PY{l+s+s1}{\PYZsq{}}\PY{l+s+s1}{Random\PYZus{}Forest}\PY{l+s+s1}{\PYZsq{}}\PY{p}{]}\PY{p}{]}\PY{o}{.}\PY{n}{mean}\PY{p}{(}\PY{n}{axis}\PY{o}{=}\PY{l+m+mi}{1}\PY{p}{)}
\PY{n}{rankings} \PY{o}{=} \PY{n}{rankings}\PY{o}{.}\PY{n}{sort\PYZus{}values}\PY{p}{(}\PY{n}{by}\PY{o}{=}\PY{l+s+s1}{\PYZsq{}}\PY{l+s+s1}{Score\PYZus{}Final}\PY{l+s+s1}{\PYZsq{}}\PY{p}{,} \PY{n}{ascending}\PY{o}{=}\PY{k+kc}{False}\PY{p}{)}

\PY{n+nb}{print}\PY{p}{(}\PY{l+s+s2}{\PYZdq{}}\PY{l+s+se}{\PYZbs{}n}\PY{l+s+s2}{\PYZhy{}\PYZhy{}\PYZhy{} TABELA COMPARATIVA DE SCORES \PYZhy{}\PYZhy{}\PYZhy{}}\PY{l+s+s2}{\PYZdq{}}\PY{p}{)}
\PY{n+nb}{print}\PY{p}{(}\PY{n}{rankings}\PY{o}{.}\PY{n}{to\PYZus{}string}\PY{p}{(}\PY{n}{index}\PY{o}{=}\PY{k+kc}{False}\PY{p}{)}\PY{p}{)}
\end{Verbatim}
\end{tcolorbox}

    \subsection{4. Pipeline de MLOps e Versionamento (Weights \&
Biases)}\label{pipeline-de-mlops-e-versionamento-weights-biases}

Nesta seção, foi implementado o versionamento de experimentos e
artefatos, uma prática mandatória para garantir a reprodutibilidade e a
governança do modelo. O uso da plataforma Weights \& Biases (W\&B)
permite o rastreamento em tempo real de cada decisão técnica tomada nas
etapas anteriores

\subsection{4.1. Governança de Dados (Data
Versioning)}\label{governanuxe7a-de-dados-data-versioning}

Não basta apenas limpar os dados; é preciso garantir que o modelo
treinado esteja vinculado à versão exata do dataset utilizado: -
Artefato raw\_data: Cópia fiel do dataset original recebido. - Artefato
clean\_data: Versão final após tratamento de outliers, imputação,
normalização e seleção via VIF/Ranking.

\subsection{4.2. Rastreamento de Hiperparâmetros e
Métricas}\label{rastreamento-de-hiperparuxe2metros-e-muxe9tricas}

Cada execução registra automaticamente o experimento: - Configuração:
Log de arquitetura (camadas, neurônios), taxa de aprendizado e dropout.
- Monitoramento de Epochs: Gráficos de convergência da função de perda
(loss) e acurácia de validação capturados a cada iteração.

\subsection{4.3. Versionamento de Modelos (Model
Registry)}\label{versionamento-de-modelos-model-registry}

O modelo final, treinado e validado, é salvo como um artefato
serializado em .pkl ou .pt, permitindo auditoria e rápida implantação

    \begin{tcolorbox}[breakable, size=fbox, boxrule=1pt, pad at break*=1mm,colback=cellbackground, colframe=cellborder]
\prompt{In}{incolor}{ }{\boxspacing}
\begin{Verbatim}[commandchars=\\\{\}]
\PY{c+c1}{\PYZsh{} Login no W\PYZam{}B}
\PY{k+kn}{import}\PY{+w}{ }\PY{n+nn}{wandb}
\PY{k+kn}{import}\PY{+w}{ }\PY{n+nn}{os}

\PY{c+c1}{\PYZsh{} Substitui \PYZsq{}TUA\PYZus{}CHAVE\PYZus{}AQUI\PYZsq{} pela chave copiada do site}
\PY{n}{os}\PY{o}{.}\PY{n}{environ}\PY{p}{[}\PY{l+s+s2}{\PYZdq{}}\PY{l+s+s2}{WANDB\PYZus{}API\PYZus{}KEY}\PY{l+s+s2}{\PYZdq{}}\PY{p}{]} \PY{o}{=} \PY{l+s+s2}{\PYZdq{}}\PY{l+s+s2}{TUA\PYZus{}CHAVE\PYZus{}AQUI}\PY{l+s+s2}{\PYZdq{}}

\PY{c+c1}{\PYZsh{} Agora o login é automático e silencioso}
\PY{n}{wandb}\PY{o}{.}\PY{n}{login}\PY{p}{(}\PY{p}{)}
\end{Verbatim}
\end{tcolorbox}

    \begin{tcolorbox}[breakable, size=fbox, boxrule=1pt, pad at break*=1mm,colback=cellbackground, colframe=cellborder]
\prompt{In}{incolor}{ }{\boxspacing}
\begin{Verbatim}[commandchars=\\\{\}]
\PY{k+kn}{import}\PY{+w}{ }\PY{n+nn}{pandas}\PY{+w}{ }\PY{k}{as}\PY{+w}{ }\PY{n+nn}{pd}
\PY{k+kn}{import}\PY{+w}{ }\PY{n+nn}{numpy}\PY{+w}{ }\PY{k}{as}\PY{+w}{ }\PY{n+nn}{np}
\PY{k+kn}{import}\PY{+w}{ }\PY{n+nn}{wandb}
\PY{k+kn}{import}\PY{+w}{ }\PY{n+nn}{joblib}
\PY{k+kn}{from}\PY{+w}{ }\PY{n+nn}{sklearn}\PY{n+nn}{.}\PY{n+nn}{neural\PYZus{}network}\PY{+w}{ }\PY{k+kn}{import} \PY{n}{MLPClassifier}
\PY{k+kn}{from}\PY{+w}{ }\PY{n+nn}{sklearn}\PY{n+nn}{.}\PY{n+nn}{model\PYZus{}selection}\PY{+w}{ }\PY{k+kn}{import} \PY{n}{train\PYZus{}test\PYZus{}split}
\PY{k+kn}{from}\PY{+w}{ }\PY{n+nn}{sklearn}\PY{n+nn}{.}\PY{n+nn}{metrics}\PY{+w}{ }\PY{k+kn}{import} \PY{n}{classification\PYZus{}report}\PY{p}{,} \PY{n}{confusion\PYZus{}matrix}\PY{p}{,} \PY{n}{f1\PYZus{}score}
\PY{k+kn}{import}\PY{+w}{ }\PY{n+nn}{matplotlib}\PY{n+nn}{.}\PY{n+nn}{pyplot}\PY{+w}{ }\PY{k}{as}\PY{+w}{ }\PY{n+nn}{plt}
\PY{k+kn}{import}\PY{+w}{ }\PY{n+nn}{seaborn}\PY{+w}{ }\PY{k}{as}\PY{+w}{ }\PY{n+nn}{sns}

\PY{c+c1}{\PYZsh{} ==========================================}
\PY{c+c1}{\PYZsh{} 1. INICIALIZAÇÃO DO W\PYZam{}B (MLOps)}
\PY{c+c1}{\PYZsh{} ==========================================}
\PY{n}{run} \PY{o}{=} \PY{n}{wandb}\PY{o}{.}\PY{n}{init}\PY{p}{(}
    \PY{n}{project}\PY{o}{=}\PY{l+s+s2}{\PYZdq{}}\PY{l+s+s2}{salary\PYZhy{}prediction\PYZhy{}mlops}\PY{l+s+s2}{\PYZdq{}}\PY{p}{,}
    \PY{n}{name}\PY{o}{=}\PY{l+s+s2}{\PYZdq{}}\PY{l+s+s2}{final\PYZhy{}model\PYZhy{}mlp}\PY{l+s+s2}{\PYZdq{}}\PY{p}{,}
    \PY{n}{config}\PY{o}{=}\PY{p}{\PYZob{}}
        \PY{l+s+s2}{\PYZdq{}}\PY{l+s+s2}{hidden\PYZus{}layer\PYZus{}sizes}\PY{l+s+s2}{\PYZdq{}}\PY{p}{:} \PY{p}{(}\PY{l+m+mi}{64}\PY{p}{,} \PY{l+m+mi}{32}\PY{p}{)}\PY{p}{,} \PY{c+c1}{\PYZsh{} Justificado: Funil para extração de features}
        \PY{l+s+s2}{\PYZdq{}}\PY{l+s+s2}{activation}\PY{l+s+s2}{\PYZdq{}}\PY{p}{:} \PY{l+s+s2}{\PYZdq{}}\PY{l+s+s2}{relu}\PY{l+s+s2}{\PYZdq{}}\PY{p}{,}          \PY{c+c1}{\PYZsh{} Evita o vanishing gradient}
        \PY{l+s+s2}{\PYZdq{}}\PY{l+s+s2}{solver}\PY{l+s+s2}{\PYZdq{}}\PY{p}{:} \PY{l+s+s2}{\PYZdq{}}\PY{l+s+s2}{adam}\PY{l+s+s2}{\PYZdq{}}\PY{p}{,}              \PY{c+c1}{\PYZsh{} Otimizador robusto para dados tabulares}
        \PY{l+s+s2}{\PYZdq{}}\PY{l+s+s2}{alpha}\PY{l+s+s2}{\PYZdq{}}\PY{p}{:} \PY{l+m+mf}{0.0001}\PY{p}{,}               \PY{c+c1}{\PYZsh{} Regularização L2}
        \PY{l+s+s2}{\PYZdq{}}\PY{l+s+s2}{learning\PYZus{}rate\PYZus{}init}\PY{l+s+s2}{\PYZdq{}}\PY{p}{:} \PY{l+m+mf}{0.001}\PY{p}{,}
        \PY{l+s+s2}{\PYZdq{}}\PY{l+s+s2}{test\PYZus{}size}\PY{l+s+s2}{\PYZdq{}}\PY{p}{:} \PY{l+m+mf}{0.2}\PY{p}{,}              \PY{c+c1}{\PYZsh{} Conforme diretriz (mínimo 20\PYZpc{})}
        \PY{l+s+s2}{\PYZdq{}}\PY{l+s+s2}{random\PYZus{}state}\PY{l+s+s2}{\PYZdq{}}\PY{p}{:} \PY{l+m+mi}{42}
    \PY{p}{\PYZcb{}}
\PY{p}{)}

\PY{c+c1}{\PYZsh{} ==========================================}
\PY{c+c1}{\PYZsh{} 2. VERSIONAMENTO DE ARTEFATOS (DATA)}
\PY{c+c1}{\PYZsh{} ==========================================}
\PY{c+c1}{\PYZsh{} Salvando o dataset limpo que você gerou na etapa anterior}
\PY{n}{X\PYZus{}final}\PY{o}{.}\PY{n}{to\PYZus{}csv}\PY{p}{(}\PY{l+s+s2}{\PYZdq{}}\PY{l+s+s2}{X\PYZus{}selected\PYZus{}features.csv}\PY{l+s+s2}{\PYZdq{}}\PY{p}{,} \PY{n}{index}\PY{o}{=}\PY{k+kc}{False}\PY{p}{)}
\PY{n}{y}\PY{o}{.}\PY{n}{to\PYZus{}csv}\PY{p}{(}\PY{l+s+s2}{\PYZdq{}}\PY{l+s+s2}{y\PYZus{}target.csv}\PY{l+s+s2}{\PYZdq{}}\PY{p}{,} \PY{n}{index}\PY{o}{=}\PY{k+kc}{False}\PY{p}{)}

\PY{n}{data\PYZus{}artifact} \PY{o}{=} \PY{n}{wandb}\PY{o}{.}\PY{n}{Artifact}\PY{p}{(}\PY{l+s+s1}{\PYZsq{}}\PY{l+s+s1}{clean\PYZus{}dataset}\PY{l+s+s1}{\PYZsq{}}\PY{p}{,} \PY{n+nb}{type}\PY{o}{=}\PY{l+s+s1}{\PYZsq{}}\PY{l+s+s1}{dataset}\PY{l+s+s1}{\PYZsq{}}\PY{p}{)}
\PY{n}{data\PYZus{}artifact}\PY{o}{.}\PY{n}{add\PYZus{}file}\PY{p}{(}\PY{l+s+s1}{\PYZsq{}}\PY{l+s+s1}{X\PYZus{}selected\PYZus{}features.csv}\PY{l+s+s1}{\PYZsq{}}\PY{p}{)}
\PY{n}{run}\PY{o}{.}\PY{n}{log\PYZus{}artifact}\PY{p}{(}\PY{n}{data\PYZus{}artifact}\PY{p}{)}

\PY{c+c1}{\PYZsh{} ==========================================}
\PY{c+c1}{\PYZsh{} 3. SPLIT ESTRATIFICADO}
\PY{c+c1}{\PYZsh{} ==========================================}
\PY{c+c1}{\PYZsh{} O stratify=y é essencial para manter a proporção da classe \PYZsq{}salary\PYZsq{}}
\PY{n}{X\PYZus{}train}\PY{p}{,} \PY{n}{X\PYZus{}test}\PY{p}{,} \PY{n}{y\PYZus{}train}\PY{p}{,} \PY{n}{y\PYZus{}test} \PY{o}{=} \PY{n}{train\PYZus{}test\PYZus{}split}\PY{p}{(}
    \PY{n}{X\PYZus{}final}\PY{p}{,} \PY{n}{y}\PY{p}{,} 
    \PY{n}{test\PYZus{}size}\PY{o}{=}\PY{n}{run}\PY{o}{.}\PY{n}{config}\PY{o}{.}\PY{n}{test\PYZus{}size}\PY{p}{,} 
    \PY{n}{stratify}\PY{o}{=}\PY{n}{y}\PY{p}{,} 
    \PY{n}{random\PYZus{}state}\PY{o}{=}\PY{n}{run}\PY{o}{.}\PY{n}{config}\PY{o}{.}\PY{n}{random\PYZus{}state}
\PY{p}{)}

\PY{c+c1}{\PYZsh{} ==========================================}
\PY{c+c1}{\PYZsh{} 4. TREINAMENTO COM EARLY STOPPING}
\PY{c+c1}{\PYZsh{} ==========================================}
\PY{n}{mlp} \PY{o}{=} \PY{n}{MLPClassifier}\PY{p}{(}
    \PY{n}{hidden\PYZus{}layer\PYZus{}sizes}\PY{o}{=}\PY{n}{run}\PY{o}{.}\PY{n}{config}\PY{o}{.}\PY{n}{hidden\PYZus{}layer\PYZus{}sizes}\PY{p}{,}
    \PY{n}{activation}\PY{o}{=}\PY{n}{run}\PY{o}{.}\PY{n}{config}\PY{o}{.}\PY{n}{activation}\PY{p}{,}
    \PY{n}{solver}\PY{o}{=}\PY{n}{run}\PY{o}{.}\PY{n}{config}\PY{o}{.}\PY{n}{solver}\PY{p}{,}
    \PY{n}{learning\PYZus{}rate\PYZus{}init}\PY{o}{=}\PY{n}{run}\PY{o}{.}\PY{n}{config}\PY{o}{.}\PY{n}{learning\PYZus{}rate\PYZus{}init}\PY{p}{,}
    \PY{n}{max\PYZus{}iter}\PY{o}{=}\PY{l+m+mi}{500}\PY{p}{,}
    \PY{n}{early\PYZus{}stopping}\PY{o}{=}\PY{k+kc}{True}\PY{p}{,}     \PY{c+c1}{\PYZsh{} MANDATÓRIO: Para evitar overfitting}
    \PY{n}{validation\PYZus{}fraction}\PY{o}{=}\PY{l+m+mf}{0.1}\PY{p}{,} \PY{c+c1}{\PYZsh{} 10\PYZpc{} do treino vira validação para o early stopping}
    \PY{n}{random\PYZus{}state}\PY{o}{=}\PY{n}{run}\PY{o}{.}\PY{n}{config}\PY{o}{.}\PY{n}{random\PYZus{}state}\PY{p}{,}
    \PY{n}{verbose}\PY{o}{=}\PY{k+kc}{True}
\PY{p}{)}

\PY{n+nb}{print}\PY{p}{(}\PY{l+s+s2}{\PYZdq{}}\PY{l+s+s2}{Iniciando treinamento da MLP...}\PY{l+s+s2}{\PYZdq{}}\PY{p}{)}
\PY{n}{mlp}\PY{o}{.}\PY{n}{fit}\PY{p}{(}\PY{n}{X\PYZus{}train}\PY{p}{,} \PY{n}{y\PYZus{}train}\PY{p}{)}

\PY{c+c1}{\PYZsh{} ==========================================}
\PY{c+c1}{\PYZsh{} 5. AVALIAÇÃO E LOG DE MÉTRICAS}
\PY{c+c1}{\PYZsh{} ==========================================}
\PY{n}{y\PYZus{}pred} \PY{o}{=} \PY{n}{mlp}\PY{o}{.}\PY{n}{predict}\PY{p}{(}\PY{n}{X\PYZus{}test}\PY{p}{)}
\PY{n}{f1} \PY{o}{=} \PY{n}{f1\PYZus{}score}\PY{p}{(}\PY{n}{y\PYZus{}test}\PY{p}{,} \PY{n}{y\PYZus{}pred}\PY{p}{,} \PY{n}{average}\PY{o}{=}\PY{l+s+s1}{\PYZsq{}}\PY{l+s+s1}{weighted}\PY{l+s+s1}{\PYZsq{}}\PY{p}{)}

\PY{c+c1}{\PYZsh{} Gerar Matriz de Confusão para o Relatório}
\PY{n}{cm} \PY{o}{=} \PY{n}{confusion\PYZus{}matrix}\PY{p}{(}\PY{n}{y\PYZus{}test}\PY{p}{,} \PY{n}{y\PYZus{}pred}\PY{p}{)}
\PY{n}{plt}\PY{o}{.}\PY{n}{figure}\PY{p}{(}\PY{n}{figsize}\PY{o}{=}\PY{p}{(}\PY{l+m+mi}{8}\PY{p}{,}\PY{l+m+mi}{6}\PY{p}{)}\PY{p}{)}
\PY{n}{sns}\PY{o}{.}\PY{n}{heatmap}\PY{p}{(}\PY{n}{cm}\PY{p}{,} \PY{n}{annot}\PY{o}{=}\PY{k+kc}{True}\PY{p}{,} \PY{n}{fmt}\PY{o}{=}\PY{l+s+s1}{\PYZsq{}}\PY{l+s+s1}{d}\PY{l+s+s1}{\PYZsq{}}\PY{p}{,} \PY{n}{cmap}\PY{o}{=}\PY{l+s+s1}{\PYZsq{}}\PY{l+s+s1}{Greens}\PY{l+s+s1}{\PYZsq{}}\PY{p}{)}
\PY{n}{plt}\PY{o}{.}\PY{n}{title}\PY{p}{(}\PY{l+s+s1}{\PYZsq{}}\PY{l+s+s1}{Matriz de Confusão \PYZhy{} Predição Salarial}\PY{l+s+s1}{\PYZsq{}}\PY{p}{)}
\PY{n}{plt}\PY{o}{.}\PY{n}{savefig}\PY{p}{(}\PY{l+s+s2}{\PYZdq{}}\PY{l+s+s2}{confusion\PYZus{}matrix.png}\PY{l+s+s2}{\PYZdq{}}\PY{p}{)}

\PY{c+c1}{\PYZsh{} Logar métricas no W\PYZam{}B}
\PY{n}{wandb}\PY{o}{.}\PY{n}{log}\PY{p}{(}\PY{p}{\PYZob{}}
    \PY{l+s+s2}{\PYZdq{}}\PY{l+s+s2}{f1\PYZus{}score\PYZus{}weighted}\PY{l+s+s2}{\PYZdq{}}\PY{p}{:} \PY{n}{f1}\PY{p}{,}
    \PY{l+s+s2}{\PYZdq{}}\PY{l+s+s2}{best\PYZus{}loss}\PY{l+s+s2}{\PYZdq{}}\PY{p}{:} \PY{n}{mlp}\PY{o}{.}\PY{n}{best\PYZus{}loss\PYZus{}}\PY{p}{,}
    \PY{l+s+s2}{\PYZdq{}}\PY{l+s+s2}{n\PYZus{}iter}\PY{l+s+s2}{\PYZdq{}}\PY{p}{:} \PY{n}{mlp}\PY{o}{.}\PY{n}{n\PYZus{}iter\PYZus{}}\PY{p}{,}
    \PY{l+s+s2}{\PYZdq{}}\PY{l+s+s2}{confusion\PYZus{}matrix}\PY{l+s+s2}{\PYZdq{}}\PY{p}{:} \PY{n}{wandb}\PY{o}{.}\PY{n}{Image}\PY{p}{(}\PY{l+s+s2}{\PYZdq{}}\PY{l+s+s2}{confusion\PYZus{}matrix.png}\PY{l+s+s2}{\PYZdq{}}\PY{p}{)}
\PY{p}{\PYZcb{}}\PY{p}{)}

\PY{n+nb}{print}\PY{p}{(}\PY{l+s+sa}{f}\PY{l+s+s2}{\PYZdq{}}\PY{l+s+se}{\PYZbs{}n}\PY{l+s+s2}{Treinamento concluído! F1\PYZhy{}Score: }\PY{l+s+si}{\PYZob{}}\PY{n}{f1}\PY{l+s+si}{:}\PY{l+s+s2}{.4f}\PY{l+s+si}{\PYZcb{}}\PY{l+s+s2}{\PYZdq{}}\PY{p}{)}
\PY{n+nb}{print}\PY{p}{(}\PY{n}{classification\PYZus{}report}\PY{p}{(}\PY{n}{y\PYZus{}test}\PY{p}{,} \PY{n}{y\PYZus{}pred}\PY{p}{)}\PY{p}{)}

\PY{c+c1}{\PYZsh{} ==========================================}
\PY{c+c1}{\PYZsh{} 6. REGISTRO DO MODELO (MODEL REGISTRY)}
\PY{c+c1}{\PYZsh{} ==========================================}
\PY{n}{model\PYZus{}filename} \PY{o}{=} \PY{l+s+s2}{\PYZdq{}}\PY{l+s+s2}{mlp\PYZus{}salary\PYZus{}model.pkl}\PY{l+s+s2}{\PYZdq{}}
\PY{n}{joblib}\PY{o}{.}\PY{n}{dump}\PY{p}{(}\PY{n}{mlp}\PY{p}{,} \PY{n}{model\PYZus{}filename}\PY{p}{)}

\PY{n}{model\PYZus{}artifact} \PY{o}{=} \PY{n}{wandb}\PY{o}{.}\PY{n}{Artifact}\PY{p}{(}\PY{l+s+s1}{\PYZsq{}}\PY{l+s+s1}{trained\PYZus{}mlp\PYZus{}model}\PY{l+s+s1}{\PYZsq{}}\PY{p}{,} \PY{n+nb}{type}\PY{o}{=}\PY{l+s+s1}{\PYZsq{}}\PY{l+s+s1}{model}\PY{l+s+s1}{\PYZsq{}}\PY{p}{)}
\PY{n}{model\PYZus{}artifact}\PY{o}{.}\PY{n}{add\PYZus{}file}\PY{p}{(}\PY{n}{model\PYZus{}filename}\PY{p}{)}
\PY{n}{run}\PY{o}{.}\PY{n}{log\PYZus{}artifact}\PY{p}{(}\PY{n}{model\PYZus{}artifact}\PY{p}{)}

\PY{c+c1}{\PYZsh{} Finalizar Run}
\PY{n}{run}\PY{o}{.}\PY{n}{finish}\PY{p}{(}\PY{p}{)}
\end{Verbatim}
\end{tcolorbox}


    % Add a bibliography block to the postdoc
    
    
    
\end{document}
